\documentclass[10pt,a4paper]{article}

% ----------------- Paquetes -------------------%
\usepackage[T1]{fontenc}
\usepackage[utf8]{inputenc}
\usepackage[a4paper,margin=2.5cm]{geometry}
\usepackage{mathptmx}             
\usepackage{changepage} 
\usepackage{setspace}
\usepackage[spanish]{babel}
\usepackage[labelfont=bf,labelsep=space]{caption}
\usepackage{amssymb}
\usepackage{amsmath}
\usepackage{ebgaramond}
\usepackage{placeins}
\usepackage{etoolbox}
\usepackage{setspace}
\usepackage{parskip} 
\usepackage{tcolorbox}
\usepackage{needspace}
\usepackage{xparse}
\usepackage{ocgx2}
\usepackage{hyperref}
\usepackage{hypcap}


%%%%%%%%%%%%%%%%%%%%%%%%%%%%%%%%%%%%%%%%%%%%%%%%%%%%%%%%%%%%%%%%
% --- Fija primero tus valores globales "buenos" ---
\setstretch{1.2}                 % Interlineado global
\setlength{\parskip}{0.7\baselineskip}
\setlength{\parindent}{0pt}
\newlength{\GlobalParskip}
\setlength{\GlobalParskip}{\parskip}

\newcounter{eqnum}[subsection]
\renewcommand{\theeqnum}{\arabic{eqnum}}
\newcounter{alphaitem}
\newcounter{numitem}

%% ------------------- Recuadros----------------------%%
\newcommand{\puntosclave}[1]{%
  \begin{tcolorbox}[
    colframe=black,
    title={Puntos clave},
    before upper={\setlength{\parindent}{0.5cm}\setlength{\parskip}{0.5\baselineskip}}
  ]
  #1
  \end{tcolorbox}
}

\newcommand{\modificaciones}[1]{
  \begin{tcolorbox}[
    colback=white,
    colframe=red,
    title={Modificaciones a realizar},
    before upper={\setlength{\parindent}{0.5cm}\setlength{\parskip}{0.5\baselineskip}}
  ]
  #1
  \end{tcolorbox}
}

%% ------------------- Listas----------------------%%
    % --- Lista alfabética ---
    \newenvironment{la}
      {\begin{list}{\alph{alphaitem})}{%
          \usecounter{alphaitem}%
          \setlength{\leftmargin}{1cm}%
          \setlength{\parskip}{3pt}% 
          \setlength{\itemsep}{0pt}% ← espacio entre ítems
          \setlength{\parsep}{0pt}%  ← párrafos dentro de un ítem
      }}
      {\end{list}}
      
    % --- Lista numérica ---
    \newenvironment{lnum}
      {\begin{list}{\arabic{numitem}.}{%
          \usecounter{numitem}%
          \setlength{\leftmargin}{1cm}%
          \setlength{\parskip}{3pt}% 
          \setlength{\itemsep}{0pt}% ← espacio entre ítems
          \setlength{\parsep}{0pt}%  ← párrafos dentro de un ítem
      }}
      {\end{list}}
      
% ------------------- Imágenes----------------------%
    \graphicspath{{Img/}}% Rutas por defecto 
    % --- Imagen adaptativa ---
    % Registros de dimensión definidos una sola vez
    \newdimen\imgcap
    \newdimen\imgmin
    \newdimen\imgmax
    \newdimen\imgremain
    \newdimen\imgh
    \newdimen\imgneed
    
    \newcommand{\imagenfit}[3]{%
      \addvspace{10pt}% separación antes
      \par\begingroup
        % Parámetros de composición (asignaciones, no \newdimen)
        \imgcap=1\baselineskip            % alto estimado de título+fuente
        \imgmin=0.15\textheight
        \imgmax=0.15\textheight
        \imgremain=\pagegoal
        \ifdim\pagegoal=\maxdimen \imgremain=0pt\fi
        \advance\imgremain by -\pagetotal
        \advance\imgremain by -\imgcap
        % Decisión de altura destino
        \ifdim\imgremain<\imgmin
          \newpage
          \imgh=0.20\textheight
        \else
          \imgh=\imgremain
          \ifdim\imgh>\imgmax \imgh=\imgmax \fi
        \fi
        % Reservar espacio para evitar cortes del bloque completo
        \imgneed=\imgh \advance\imgneed by \imgcap
        \Needspace{\imgneed}
        % Cuerpo
        \noindent\begin{minipage}{\textwidth}
          \centering
          \captionof{figure}{\textemdash\ #2}\par\vspace{0.3em}% título arriba
          \includegraphics[width=\linewidth,height=\imgh,keepaspectratio]{\detokenize{#1}}\par
          \vspace{0.3em}\small\textit{Fuente: }#3% fuente abajo
        \end{minipage}%
      \endgroup
      \par\addvspace{10pt}%
    }

%------------ Bloque de RECUADRO ESPECIAL -------------%
    \newenvironment{specialblock}[1]{%
      \begin{quote} % crea sangría a izquierda y derecha
      \small % texto más pequeño
      \noindent\textbf{#1}\\[-0.3em] % título en negrita
      \rule{\linewidth}{0.4pt}\\[0.6em]}{
      \end{quote}}
      

%-------------------------- Author Font -----------------------%
    \newcommand{\authorfont}[1]{{\fontfamily{EBGaramond-TLF}\selectfont #1}}

%-------------------------- Anexos -----------------------%
    \makeatletter
    \newcounter{anexo}
    \renewcommand{\theanexo}{\Roman{anexo}}
    \newcommand{\anexo}[1]{%
      \clearpage
      \phantomsection
      \refstepcounter{anexo}%
      \noindent\textbf{Anexo \theanexo.- }#1\par\medskip
      \addcontentsline{stc}{section}{Anexo \theanexo.- #1}%
    }
    \makeatother

    %-------------------------- Lista de figuras (personal) -----------------------%
\makeatletter
\newcommand{\MyListOfFigures}{%
  \clearpage
  \phantomsection
  {\centering\bfseries\Large Índice de figuras\par}%
  \medskip
  % Entrada en nuestro índice de contenido (stc)
  \addcontentsline{stc}{section}{Índice de figuras}%
  % Recuperación del listado (sin cabecera propia)
  \begingroup
    \parindent=0pt\parskip=2pt
    \@starttoc{lof}%
  \endgroup
}
\makeatother

%-------------------------- Bibliografía -----------------------%
\makeatletter
% Contador para las referencias
\newcounter{myref}
% Cabecera de la bibliografía, entra en el índice 'stc'
\newcommand{\MyBibliography}{%
  \clearpage
  \phantomsection
  {\centering\bfseries\Large Bibliografía y referencias\par}%
  \medskip
  \addcontentsline{stc}{section}{Bibliografía y referencias}%
  \setcounter{myref}{0}%
}
\newcommand{\MyBibItem}[4]{%
  \refstepcounter{myref}%
  \noindent[\themyref]\ #1.\ \emph{#2}. #3. #4\par
}
\makeatother

%--------- Establecer cuatro niveles de índice --------%
    \setcounter{tocdepth}{4}   
    \setcounter{secnumdepth}{4}% numera hasta \paragraph

%%%%%%%%%%%%%%%%%%%%%%%%%%%%%%%%

\makeatletter

    %--------- Interlineado Global --------%
    \edef\GlobalBaseLineStretch{\baselinestretch}
    
    %--------- Espaciado de seccinones --------%
    \renewcommand\section{\@startsection{section}
      {1}{\z@}
      {2.5ex \@plus 1ex \@minus .2ex} % espacio antes
      {1.5ex \@plus .2ex}             % espacio después
      {\normalfont\Large\bfseries}    % formato del título
    }
    \renewcommand\subsection{\@startsection{subsection}{2}{\z@}
      {3ex \@plus 0.8ex \@minus .2ex}
      {1.5ex \@plus .2ex}
      {\normalfont\large\bfseries}
    }
    \renewcommand\subsubsection{\@startsection{subsubsection}{3}{\z@}
      {3ex \@plus 0.5ex \@minus .2ex}
      {1ex \@plus .2ex}
      {\normalfont\normalsize\bfseries}
    }
    \renewcommand\paragraph{\@startsection{paragraph}{4}{\z@}%
    {-2ex \@plus -0.5ex \@minus -.2ex}% espacio antes
    {0.8ex \@plus .2ex}% espacio después
    {\normalfont\normalsize\itshape}}% ← cursiva en lugar de negrita

    %--------- Ecuaciones --------%   
    % Variables globales para almacenar títulos y notas
    \gdef\leq@current@section{}%
    \gdef\leq@current@subsection{}%
    \gdef\leq@last@printed@section{}%
    \gdef\leq@last@printed@subsection{}%
    
    % Comando para añadir nota (invisible en el cuerpo)
    \newcommand{\eqnote}[1]{%
      \addcontentsline{leq}{leqnote}{#1}%
    }
    
    % Comando para ecuaciones
    \newcommand{\eqblock}[2]{%
      % Verificar si necesitamos imprimir el título de sección
      \ifx\leq@current@section\leq@last@printed@section\else
        \addcontentsline{leq}{leqsection}{\leq@current@section}%
        \global\let\leq@last@printed@section\leq@current@section
        \gdef\leq@last@printed@subsection{}% Reset subsección al cambiar sección
      \fi
      % Verificar si necesitamos imprimir el título de subsección
      \ifx\leq@current@subsection\leq@last@printed@subsection\else
        \ifx\leq@current@subsection\@empty\else
          \addcontentsline{leq}{leqsubsection}{\leq@current@subsection}%
          \global\let\leq@last@printed@subsection\leq@current@subsection
        \fi
      \fi
      % Ecuación
      \refstepcounter{eqnum}%
      \begin{equation}
        #1\tag{E.\theeqnum}
      \end{equation}%
      \addcontentsline{leq}{leqt}{[E.\theeqnum] -- #2}%
      \addcontentsline{leq}{leqm}{#1}%
    }
    
    %%% ----- Índice de ecuaciones ----------%%%
    % Tamaño para []
    \newcommand{\leqIndexFont}{\normalsize}
    \newcommand{\leqTitleFont}{\tiny}
    \newcommand{\leqNoteFont}{\tiny}
    % Cabecera de sección 
    \newcommand\l@leqsection[2]{%
      \addvspace{25pt}% Mayor separación antes de nueva sección
      \noindent\leqIndexFont
      \makebox[\linewidth]{\rule{.38\linewidth}{0.4pt}\ \ \textbf{  \MakeUppercase{#1}}\ \ \rule{.38\linewidth}{0.4pt}}%
      \par\vspace{-10pt}
    }
    % Cabecera de subsección 
    \newcommand\l@leqsubsection[2]{%
      \par\addvspace{15pt}%
      \noindent\leqIndexFont #1
      \par\vspace{-7pt}
    }
    % Título de ecuación 
    \newcommand\l@leqt[2]{%
    \par\addvspace{10pt}%
      \par\noindent\leqTitleFont #1\par
    }
    % Miniatura de ecuación
    \newcommand\l@leqm[2]{%
      \vspace{0.3\baselineskip}%
      \par\scriptsize\noindent\hspace{1.5em}\(#1\)\par%
    }
    % Nota de ecuación 
    \newcommand\l@leqnote[2]{%
      \vspace{0.2\baselineskip}%
      \par\noindent\leqNoteFont\hspace{1.5em}\textbf{Nota.--} #1\par\vspace{0.5\baselineskip}%
    }
    % Generación del índice
    \newcommand{\mylistofequations}{%
      \clearpage
      \phantomsection
      % Título visual de la lista (ajústalo a tu gusto):
      {\centering\bfseries\Large Índice de ecuaciones\par}
      % Entrada en nuestro índice 'stc':
      \addcontentsline{stc}{section}{Índice de ecuaciones}%
      % Llamada a tu lista real (usa el nombre exacto de tu macro):
      \@starttoc{leq}%
    }
    
    %--------- Captura de títulos de sección --------%
    \let\leq@original@section\section
    \let\leq@original@subsection\subsection
    % Flag para saber si estamos en una sección asterisco
    \newif\ifLeqStarSection
    \LeqStarSectionfalse
    
    % Redefinir \section CON CONTROL DE ASTERISCO
    \renewcommand{\section}{%
      \@ifstar{\leq@section@star}{\@dblarg\leq@section@nostar}%
    }
    \def\leq@section@star#1{%
      \global\LeqStarSectiontrue
      \gdef\leq@current@section{#1}%
      \gdef\leq@current@subsection{}%
      \leq@original@section{#1}%
      \global\LeqStarSectionfalse
    }
    \def\leq@section@nostar[#1]#2{%
      \global\LeqStarSectionfalse
      \gdef\leq@current@section{#2}%
      \gdef\leq@current@subsection{}%
      \leq@original@section[#1]{#2}%
    }
    % Redefinir \subsection
    \renewcommand{\subsection}{%
      \@ifstar{\leq@subsection@star}{\@dblarg\leq@subsection@nostar}%
    }
    \def\leq@subsection@star#1{%
      \global\LeqStarSectiontrue
      \gdef\leq@current@subsection{#1}%
      \leq@original@subsection{#1}%
      \global\LeqStarSectionfalse
    }
    \def\leq@subsection@nostar[#1]#2{%
      \global\LeqStarSectionfalse
      \gdef\leq@current@subsection{#2}%
      \leq@original@subsection[#1]{#2}%
    }

    % ----- Restauración de valores Globales de estilo ----------%
    % Flags
    \newif\ifInFootnote
    \InFootnotefalse
    \pretocmd{\@footnotetext}{\InFootnotetrue}{}{}
    \apptocmd{\@footnotetext}{\InFootnotefalse}{}{}
    % Profundidad de listas (soporta anidadas)
    \newcount\MyListDepth \MyListDepth=0
    \AtBeginEnvironment{la}{\advance\MyListDepth by 1}
    \AfterEndEnvironment{la}{\advance\MyListDepth by -1}
    \AtBeginEnvironment{lnum}{\advance\MyListDepth by 1}
    \AfterEndEnvironment{lnum}{\advance\MyListDepth by -1}
    \AtBeginEnvironment{itemize}{\advance\MyListDepth by 1}
    \AfterEndEnvironment{itemize}{\advance\MyListDepth by -1}
    \AtBeginEnvironment{enumerate}{\advance\MyListDepth by 1}
    \AfterEndEnvironment{enumerate}{\advance\MyListDepth by -1}
    % Restauración diferida: hazlo al INICIO del próximo párrafo real
    \newif\if@pendingRestore
    \@pendingRestorefalse
    \newcommand{\RequestRestore}{\global\@pendingRestoretrue}
    \everypar{%
      \if@pendingRestore
        \global\@pendingRestorefalse
        \ifInFootnote\else
          \ifnum\MyListDepth=0
            \setlength{\parskip}{\GlobalParskip}%
            \setstretch{\GlobalBaseLineStretch}%
          \fi
        \fi
      \fi
    }
    % Al final de bloques:
    \AfterEndEnvironment{equation}{\RequestRestore}
    \AfterEndEnvironment{equation}{\RequestRestore}
    \AfterEndEnvironment{la}{\RequestRestore}
    \AfterEndEnvironment{lnum}{\RequestRestore}
    % Al final de macros:
    \apptocmd{\eqblock}{\RequestRestore}{}{}
    \apptocmd{\imagenfit}{\RequestRestore}{}{}
    
\makeatother

% ===================== ToC propio con 4 niveles (único bloque) =====================
\makeatletter

% --- Escritores: añadimos una línea al .stc en cada cabecera numerada ---
% Guardamos las versiones actuales para llamar al comportamiento normal
\let\MyTOC@orig@section\section
\let\MyTOC@orig@subsection\subsection
\let\MyTOC@orig@subsubsection\subsubsection
\let\MyTOC@orig@paragraph\paragraph

% SECTION
\renewcommand{\section}{%
  \@ifstar{\MyTOC@section@star}{\@dblarg\MyTOC@section@nostar}%
}
\def\MyTOC@section@star#1{%
  \MyTOC@orig@section{#1}% sin entrada al .stc
}
\def\MyTOC@section@nostar[#1]#2{%
  \MyTOC@orig@section[{#1}]{#2}% comportamiento normal (escribe al .toc)
  % Escribimos también nuestra línea al .stc (texto con \numberline)
  \addcontentsline{stc}{section}{\protect\numberline{\thesection}#2}%
}

% SUBSECTION
\renewcommand{\subsection}{%
  \@ifstar{\MyTOC@subsection@star}{\@dblarg\MyTOC@subsection@nostar}%
}
\def\MyTOC@subsection@star#1{%
  \MyTOC@orig@subsection{#1}%
}
\def\MyTOC@subsection@nostar[#1]#2{%
  \MyTOC@orig@subsection[{#1}]{#2}%
  \addcontentsline{stc}{subsection}{\protect\numberline{\thesubsection}#2}%
}

% SUBSUBSECTION
\renewcommand{\subsubsection}{%
  \@ifstar{\MyTOC@subsubsection@star}{\@dblarg\MyTOC@subsubsection@nostar}%
}
\def\MyTOC@subsubsection@star#1{%
  \MyTOC@orig@subsubsection{#1}%
}
\def\MyTOC@subsubsection@nostar[#1]#2{%
  \MyTOC@orig@subsubsection[{#1}]{#2}%
  \addcontentsline{stc}{subsubsection}{\protect\numberline{\thesubsubsection}#2}%
}

% PARAGRAPH
\renewcommand{\paragraph}{%
  \@ifstar{\MyTOC@paragraph@star}{\@dblarg\MyTOC@paragraph@nostar}%
}
\def\MyTOC@paragraph@star#1{%
  \MyTOC@orig@paragraph{#1}%
}
\def\MyTOC@paragraph@nostar[#1]#2{%
  \MyTOC@orig@paragraph[{#1}]{#2}%
  % Si no numeras los \paragraph, cambia \theparagraph por texto plano sin \numberline
  \addcontentsline{stc}{paragraph}{\protect\numberline{\theparagraph}#2}%
}

% --- Impresor del índice propio (lee el .stc) ---
\newcommand{\MyTOC}{%
  \par\bigskip
  {\centering\bfseries\Large Índice de contenido\par}%
  \medskip
  \begingroup
    \parindent=0pt\parskip=2pt
    \@starttoc{stc}%
  \endgroup
}

\makeatother
% ===================== fin ToC propio =====================


%%%%%%%%%%%%%%


% --- Datos del tema ---
\title{Política Presupuestaria (III): Finanzas de las Administraciones Públicas en España\\
\large Contabilidad Nacional. Situación actual y perspectivas según el PE. Estructura de ingresos y presión fiscal. Estructura del gasto por funciones según la clasificación COFOG. Distribución de gastos e ingresos públicos por subsectores de las Administraciones Públicas}
\author{Marcelo Ortega}
\date{Septiembre 2026}

\usepackage{soul}\sethlcolor{yellow}% resaltado amarillo: \hl{texto} (Panel TCEE, ⌃H)
\begin{document}
\maketitle
\newpage

% --- Página de resumen y modificaciones ---
\puntosclave{ %% Escribir puntos clave Apuntes CECO
     
    }
\vspace{1cm}

\modificaciones{
    
    }

\newpage
\MyTOC
\newpage

\mylistofequations
\clearpage

\MyListOfFigures
\clearpage


\pagenumbering{arabic}
\setcounter{page}{1}


\section{Introducción}



\subsection{Contextualización}


\subsubsection{Relevancia}




\subsubsection{Problemática}



\subsection{Estructura de la exposición}




Este tema tiene por objetivo conocer el \textbf{contenido de los ingresos y los gastos de las Administraciones Públicas} españolas.

El tema tiene un perfil memorístico innegable, así como la conveniencia de actualizar sus cifras periódicamente. Por lo que se refiere al primer inconveniente. las cifras se proporcionan en términos de PIB. lo cual hace el tema más lle,adcro: además de poner de relevancia la importancia macroeconómica de las magnitudes. En cuanto al segundo inconveniente, no se espera que los tribunales vayan a exigir cifras exactas y actualizadas, sino que el opositor conozca el orden de magnitud de ciertos agregados importantes: déficit público (o superávit).\footnote{%
    Quizás el dato de saldo público (como porcentaje del PIB) del año anterior y el esperado por el gobierno para el año en curso, sí deban darse con exactitud. pues serán generalmente conocidos.
} gasto total, ingreso total, gasto en prestaciones sociales, inversión pública, e ingresos por impuestos directos o indirectos. 

No obstante, en aras de una mayor claridad explicativa, en el tema figuran cifras exactas, redondeadas a una décima del PIB. o a dos décimas en el caso de que interese que el lector pueda comprobar la suma/resta de magnitudes (como en los casos de las operaciones financieras y del desglose de los niveles de las AAPP). En la exposición del tema, el opositor deberá adaptar el lenguaje en función de su capacidad memorística y del tiempo de que disponga para actualizar cifras. En general, debe bastar con decir que se sitúa \textit{en torno a .. o ·'por encima de .. } un determinado porcentaje del PIB, o que representa un determinado porcentaje del total.

Las cifras que se consignan en este tema son las vigentes en primavera de 2023. Comprenden las cuentas de las AAPP liquidadas por la IGA E hasta 2022 y las de la Actualización del Programa de Estabilidad 2023-2026. El 18 de septiembre de 2023. el INE publicó una revisión de la contabilidad nacional anual que corrigió levemente al alza el Pll3 en 2020. 2021 y 2022. Esta revisión no se ha incorporado al tema, por estar éste ya redactado y por no cambiar esta revisión las principales ideas del mismo. No obstante, por si es de interés para el opositor. en el Anexo 4 del tema se incluye la tabla 1 con las ratios de ingresos y gastos tal como quedan después de esta revisión. En otoño de 2023 la lGAE publicará un mayor ele desglose de ingresos. pero no es muy relevante para los propósitos del tema. Las grandes cifras de ingresos y gastos len euros) de esta publicación no suelen cambiar con respecto a las publicadas en primavera. lgualmente, el 15 de octubre de 2023, se publicará el Plan Presupuestario del Reino de España. con nuevas previsiones rnacroeconómicas y fiscales para 2023 y 2024.

El 3 1 de marzo de 2024 la IGAE, en la Contabilidad Nacional Trimestral, publicará las cuentas del cuarto trimestre <le 20232 y. por tanto. todo el año 2023\footnote{%
    Todos los 30 o 31 de mar.lo de cada año se publica la CNTR con el cuarto trimestre del año anterior. con las principales macro magnitudes. incluyendo las de la�� AAPP-
} cuyos datos en este tema figuran como previsión de programa de Estabilidad 2023-2026. Así mismo. el 30 de abril ele 2024 sepublicará la Actualización del Programa de Estabilidad 2024-2027. Estas dos publicaciones contendrán los datos vigentes cuando haya que defender este tema en la OEP de 2023. En el caso de que el opositor quiera actualizar las cifras. puede seguir la estructura del tema y sustituir con cierta facilidad los datos con la información que publicará la IGAE y el Programa de Estabilidad, que año tras año siguen el mismo orden de presentación. 

La principal fuente bibliográfica es el Informe trimestral de la IGAE (2023 a). Allí el opositor e11contrar:\ las cuentas de las distintas AAPP con el desglose que aparece en este tema y en el Programa de Estabilidad. Si se quiere un detalle mayor, hay que recun-ir a los datos de IGAE (2023b). Por ejemplo, en el apartado II se dividen los ingresos y gastos entre corrientes y de capital, pues esta división tiene alguna relevancia económica. Sin embargo, esta división no aparece ni en IGAE (2023a) ni en el Programa de Estabilidad. Si el opositor quiere mantenerla al actualizar los datos. debe recmTir a IGAE (2023b).

En el caso de algunas fuentes, como el Boletín Estadístico del Banco de España, se consigna el número del cuadro de donde proceden los datos para facilitar la actualización de los mismos por parte del opositor que lo desee.

Como se comprobará, el tema contiene una gran cantidad de tablas. No son para memorizar, sino para información del opositor y para que escoja comentar los datos que prefiera, además o en lugar de los que se comentan en el texto.

Por último y en relación con todo lo anterior. si bien el enfoque adoptado en este documento es esencialmente descriptivo (con la finalidad, por un lado, de estructurar el amplísimo contenido y, por otro, de ser relativamente exhaustivo al respecto siguiendo el título del BOE y lo que puede pedir el Tribunal), la exposición oral se enriquecerá sustancialmente de introducir el opositor valoraciones analíticas. Es decir. es co11ditio sine qua 11011 conocer la información ·'objetiva•·, pero el valor añadido que puede aportar el opositor consiste en basarse en ésta para realizar auténtico análisis. Para este fin. este terna: (i) introduce reflexiones básicas y ampliamente aceptadas, (ii) señala numerosas fuentes bibliográficas con artículos analíticos que incorporarán cierta subjetividad en los resultados. Además, el opositor debe ser capaz de ligar ( evitando solapamientos) este tema con los otros acerca del Presupuesto español - 4B21 a 4826-, así corno con el tema 4B I sobre contabilidad del sector público.


------------------------------------ -> Introducción

Este tema tiene por objeto conocer el tamaño del sector público dentro de la economía. la composición de sus ingresos y gastos ) . en este último caso, su estructura por funciones. Los datos a los que nos referiremos son los prororcionados por la contabilidad nacional. según la metodología del SEC-20 I O\ en el que, dentro de los Sectores Institucionales (Sociedades no Financieras. Instituciones Fi nancieras. ! logares y Admini straciones Públicas (AA PP)). las AAPP constituyen el Sector S. 1 3 . Las operaciones de las AA PP se registran. según corresponda. en las cuentas de Producción. de Distribución y Utilización de Renta. y de Acumulación. A estas cuentas hay que aiiadir los balances. en los que se consignan los activos ) pasivos de las A A PP.

El procedimiento de contabilización de las operaciones de las AAPP queda fuera del objeto del tema. el cual se centra en una presentación resumida de las finanzas públicas en unas cuentas que recogen las principales rúbricas de Recursos ( I ngresos) ) Empleos (Gastos). En el caso de los balances. nos referiremos únicamente a la deuda de las A A PP.

Para ello. cornenzaremos considerando el conjunto de las Administraciones Públicas (A A PP) a través de su cuenta consolidada. con J¡¡s ci fras en términos de contabilidad nacional. como se ha dicho. Se trata de una visión más simple y resumida en un país con una Admin istración muy descentralizada. pero que es muy úti l desde el punto de vista económ ico por cuanto los ci udadanos pagan unos tributos y tasas y obtienen a cambio bienes, servicios y prestaciones; no siendo en este enfoque importante qué administración recibe estos i ngresos ni qué adm inistración proporciona estos bienes, servicios y prcstaciones.\footnote{%
    La cuenta consolidada de las AAPP es. en definiti, a. una cuenta bastante resumida (ver Tabla l. más adelante) en la que se recogen wdo�� los ingre��os que las distintas AAPP perciben de los demás Sectores Institucionales y del Resto del Mundo, y todos lo:. gastos que las AA PP realizan. sin desglosar los subsectores de la Administración que son destino u origen de esos ingresos y gastos.
} Además, a nivel macroeconómico y de supervisión macro-fiscal ¿5te es el enfoque que interesa. Cerraremos el estudio de la cuenta consolidada de las A A P P con la clasificación funcional del gasto y con el análisis de la presión fiscal en España.

Después de este enfoque consolidado, se exam inan las finanzas de los distintos subsectores de las A A PP: Administración Central. Administración Regional. Adm inistración Local y Fondos de la Seguridad Social.

Las magnitudes se proporcionan en porcentaje del PIB, pues reflejan mejor la importancia de cada rúbrica. El año al que se referirán los datos es 2022 ( último con datos publ icados al finalizar la redacción del tema) y se harán comparaciones con años anteriores, en especial con 20 1 9 ( ai'io anterior a la crisis de la Covid-19), as í como con las perspectivas de futuro según el Progrnma de Estabilidad 2023-2026.

Por s u mayor familiaridad. la terminología que se emplea en este tema es la misma que la del Programa de Estabilidad de Espal'ía. que es distinta de la empleada por la I ntervención General de la Administración del Estado (IGAE) en sus cuentas públicas en términos de c ontabilidad nacional. Así. por ejemplo. emplearemos los términos '•ingresos·· y '·gastos"' en vez de ·'recursos'" y ·'empleos'·: al igual que '•impuestos directos"' e "'impuestos i ndirectos '" en lugar de '"impuestos sobre la renta y el patrimonio'· e •'impuestos sobre la producción y las importaciones·'. Así mismo. en ocasiones emplearemos los términos ·'superávit'' y '•déficit'", en vez de capacidad o neces idad de linanciación.




\clearpage
\section{Cuenta Conslidada}
\puntosclave{
    La crisis de la Co\ id-19 ha intcrrumrido el proce::-o de ajuste del déficit de las AAPP que se venía produciendo desde 2012. Tras la crisis. las AA PP parecen instalada!> en unos niveles de ingresos y gastos. como porcentaje del PIB, mayores que antes (2019)
    
    • F.n 2022 los gastos no financieros totales ascendieron al 4 7 .8 por ciento del PI B, 5,.5 puntos má!> que en 2019. Se espera que esta ratio baje algo en el medio plazo. pero permanecerá alta.

    • Entre 2008 y 2023 los gastos corrientes per cápita habrían crecido un 26.4 por ciento, m icntras que el Plíl per cápita lo habría hecho en un 3 por ciento. Esta divergencia de tasas representa un 1iesgo para la sostenibilidad del estado del bienestar.

    • Los ingresos no financieros totales ascendieron en 2022 al 43 por ciento dd PIB. 3.8 puntos por encima de 20 19. En cambio. los in1=,rresos netos de transferencias y sub\'enc1ones fueron el 1 7.4 por ciento del PIB. 1.6 puntos menos que en 2019.
    
    • En 2022 la necesidad de financiación (déficit) fue del 4.8 por ciento del Pll3. Otros saldos fueron también negativos: el ahorro público (-2 por ciento del PIBJ. el saldo público primario (-2.4 por ciento del PI B) y el saldo público de la Regla de Oro (-0,8 por ciento del PIB).

    • Se espera que estos desequilibrios se corrijan en el medio plazo, continuando la línea anterior a la crisis de la Covid-19. Además. la activación de las reglas macro- fiscales en la VE puede obligar a un ajuste mayor.
    
    • Ariadiendo a la necesidad de financiación de las AAPP (4.8 por ciento del PIB ) la variación neta <le activos financieros y unos ajustes. resulta en 2022 una variación. en euros, de la deuda PDE equivalente al 5.7 por ciento del PlB.
    
    • La principal rúbrica de gasto son las transferencias sociales. que suponen el 20 .1 por ciento del PIB. De ellas. las más importantes son las Prestaciones Sociales No en Especie. con el 17.2 por ciento del PIB. Son el concepto que más crece (un 67 por ciento desde 2008) y más se espera que crezca. por el crecimiento del número de pensiones. su revalorización anual y el mayor valor de las nuevas pensiones.
    
    • La remuneración de asalariados y los consumos intermedio�� suman un gasto equivalente al 17.5 por ciento del PIB. Mientras que la Formación Bruta de Capital asciende al 2 .8 por ciento del PIB. Por lo tanto, la demanda agregada del sector público es del 20.3 por ciento del PIB.

    • La principal rúbrica de ingresos son las Cotizaciones Sociales, con el 13.6 por ciento del PIB. 

    • La suma de impuestos directos e indirectos asciende al 24.5 por ciento del PIB. con un peso un poco mayor de los impuestos directos.

    • Los datos contenidos en la cuenta de las AAPP, non11alrnente descendiendo a cierto nivel de detalle, son los que se emplean para evaluar el cumplimiento de las reglas fiscales dedéficit y gasto.
    • A las operaciones no financieras recogidas en la cuenta de las AAPP hay que añadirles las operaciones financieras, para obtener la necesidad de endeudamiento de las AAPP. Sobre esta cuestión se inserta una nota en el Recuadro 1.
}

Empezamos examinando las principales rúbricas de operaciones no financieras de ingresos y gastos del sector AAPP. Los datos que son previsiones en este tema (de 2023 en adelante), no reflejan los flujos relacionados con el MR.R y el PRTR.5. Esto es así por la nonna de Eurostat que establece la neutralidad, a efectos del cálculo del déficit público, de los ingresos y gastos públicos de esta iniciativa cuando los ingresos recibidos por las AAPP nacionales desde la UE sean transferencias6. Al finalizar el examen de los principales resultados. y antes de ver en detalle los ingresos y los gastos no financieros, nos referiremos brevemente a las operaciones financieras. 

Desde 2012. año en el que las finanzas públicas se resintieron más por efecto de la Gran Recesión. la economía española fue corrigiendo gradualmente sus desequilibrios principales: fiscal, exterior, desempleo etc. En el ámbito fiscal. las AAPP pasaron de una necesidad de financiación en 201 2 del 10,7 por ciento del PIB. al 3, 1 por ciento del PlB en 2019. Sin embargo, la crisis de la Covid- 19 supuso un aumento del gasto público y un hundimiento del PlB nominal. lo que llevó a la ratio ingresos/PIB a superar el 50 por ciento por primera vez en la historia. Los ingresos públicos tam bién cayeron, pero en menor proporción que el PJB nominal. de ahí que la ratio sobre PIB también subiera. aunque menos que en el caso del gasto. Esto llevó a un aumento de la Necesidad de Financiación de las AAPP hasta el 10.1 por ciento del PIB en 2020. En este año y en 2021 fueron muy numerosas las mediadas de gasto (al alza) y de ingresos (al alza y a la baja) adoptadas. tanto transitorias como pe1manentes, para combatir la crisis. Las perspectivas a mediados de 2023 apuntan a que las AAPP se instalarán en el futuro en unos niveles de gastos e ingresos. como porcentaje del PIB, sustancialmente mayores que antes de la crisis. A continuación. veremos esto en detalle.

En la tabla l se ofrece el detalle de la evolución de las principales rúbricas de ingresos y gastos no financieros para el periodo 2019-2022. con las previsiones para 2023 y 2024 del Programa de Estabilidad de España 2023-20267


\textbf{TABLA RESUMEN: ADMINISTRACIONES PÚBLICAS (USOS Y FUENTES) CONSOLIDADO}

11.l. Principa les resultados

Ga.,·tos totales
 
En 2022. los gastos no financieros totales de las AAPP ascendieron al 47,8 por ciento del PIB; es decir, 5,5 puntos por encima del nivel de 20 1 9, último año antes de la crisis de la Covid- 1 9 ( 42,3 por ciento del PIB). Las previsiones a mediados de 2023 apuntaban al mantenimiemo del gasto público a unos niveles algo más bajos en el medio plazo, pero sin bajar hasta el nivel de 20 1 9. En los años de la crisis de la Covid- 1 9 (2020 y 202 1 ) el gasto público aumentó y el PIB cayó bastante por debajo del nivel de 201 9, con l a consiguiente elevación de la ratio con respecto al PIB nominal.

Esta cifra nos proporciona una primera idea del tamaño del sector público en España. En ese mismo año, ese gasto se situó en el 50,8 por ciento del PIB en el conjunto de la UE. en el 46,8 por ciento en d Reino Unido y en el 38,4 por ciento en Estados Unidos.\footnote{%
    Para comparaciones de gasto con la UE. incluyendo la clasificación funcional (COFOG), ver Alloza et uf. (2022)
}

Ciastos corrientes y sostenibilidad.fiscal

Los gastos corrientes fueron en 2022 el 9 1 .7 por ciento de los gastos totales. Llama la atención el rápido crecimiento de esta magnitud en términos per cápita. comparado con el crecimiento mucho más lento del PIB per cápita. Así, por ejemplo, de 2008 a 2023 el gasto corriente per cápita a precios constantes habría aumentado un 26,4 por ciento, mientras que el PIB real per cápita lo habría hecho en un 3,0 por ciento.

Esto plantea un problema de sostenibilidad del estado del bienestar que, como veremos, está detrás de la mayor parte del gasto público corriente. Si el gasto se acaba financiando con impuestos, la carga fiscal sobre los contribuyentes habrá aumentado más que su renta. Si esta tendencia se prolonga, las alternativas serían, o rediseñar las políticas de gasto o conseguir un mayor crecimiento económico. Pero esto último es lo que intenta un gobierno en todo caso, ingresos totales

Por su parte, los ingresos no financieros totales de las AAPP ascendieron al 43,0 por ciento del PIB en ese m ismo año. 3,8 puntos más que en '.!0 1 9 (39,2 por ciento del PfB). Esta ratio se situó en la U E en el 46,5 por ciento, en e l Reino Unido en e l 4 1 , l por ciento y en Estados Unidos en e l 34,4 por ciento.

Ingresos públicos netos de tra11sfere11ciasy subvenciones

La ratio ingresos públicos no financieros como porcentaje del PIB, 43,0 por ciento en 2022, nos da una idea de la carga del sector público sobre los contribuyentes. pero, para este propósito, es más acertado restar de los ingresos públicos todos los gastos de transferencias (corrientes y de capital) y de subvenciones que la Administración "devuelve'· al sector privado (y al sector exterior). Si hacernos esto, obtenemos una ratio de ingresos públicos netos de transferencias y subvenciones del 1 7,4 por ciento del PIB en 2022.\footnote{%
    También se puede hacer lo mismo para obtener el gasto público neto de transferencias y subvenciones (que equivale a la suma del consumo público. la inversión pública y los gastos de intereses). Es. quizás. una medida mejor del peso del sec1or público, dado que el gasto público a la larga hay que pagarlo en su totalidad con impuestos, aunque a corto plazo se recurra al endeudamiento si los ingresos son insuficientes y hay déficit.
} Esta cifra es inferior al 1 9 por ciento del PlB que se registró en 20 1 9. Más adelante abordaremos con más detalle el tema de la carga del sector público cuando nos refiramos a la presión fiscal.

Necesidad de Ji11u11ciaciú11

La necesidad de financiación de las AAPP ascendió al 4.8 por ciento del PIB. mientras que en la curozona fue el 3,6 por ciento del PIB. en la UE el 3.4 por e1ento del PIB, t'n el Reino Unido el 5,2 por ciento del PIB y en Estados Unidos el 4.0 por ciento del PIB.e espera que la necesidad de financiación descienda en 2023 al 3.9 por ciclllo del PIB ) que 2024 se sitúe en el 3.0 por ciento. En 2025 y 2026 (no consignados en la tabla 1 ) se proseguiría el ajuste fiscal hasta una necesidad de financiación del 2.5 por ciento del PI B. La activación de las reglas macro- fiscales en la UE puede llevar a un ajuste mayor.

Ahorro público

Además de la necesidad de financiación. otros saldos tienen también interés. Así, por ejemplo. el ahorro público; es decir. la diferencia entre ingresos corrientes y gastos corrientes. ascendió en 2022 al -2,0 por ciento del PlB

Si al ahorro público le sumamos los impuestos sobre el capital y el saldo neto de transferencias de capital (ingresos de capital menos gastos de capital). nos queda el excedente que sirve para financiar la inversión pública (formalmente llamada Formación Bruta de Capital público). Este excedente no fue tal porque tuvo signo negativo: el - 1 .9 por ciento del PIB.\footnote{%
    Excedente = Ahorro público + Impuestos de capital +(Ingresos por transferencias de capital - Gastos por tramsferencias de capital) = -2,0 + 0,4 + (0,8 - 1 ,1) = - 1,9
} Como la inversión pública ascendió al 2,8 por ciento del PIB, el déficit (añadiendo la rúbrica de adquisiciones menos cesiones de activos no producidos) fue el 4.8 por ciento del PIB anteriormente reseñado.

La información proporcionada en la cuenta de las AAPP del Programa de Estabilidad no tiene un desglose que permite saber la evolución futura del ahorro público. pero del ajuste total previsto del déficit público a medio plazo se puede inferir que este desequilibrio se corregirá. 

Saldo primario

El saldo público primario; es decir. la diferencia entre ingresos totales y los gastos totales sin contar con el gasto de intereses de la deuda. ascendió al -2,;4 por ciento del PIB. Según las proyecciones del Programa de Estabilidad 2023-2026, en 2025 el saldo primario sería nulo y habría un pequeño superávit primario en 2026 (0,1 por ciento del PIB).

Regla de Oro

La necesidad de financiación. neta de los gastos de capital, ascendió al -0,8 por ciento del PIB. Según la llamada Regla de Oro el déficit público no debe ser mayor que los gastos de capital; es decir. un mayor endeudamiento está justificado si es para financiar gastos de capital. pues estos ayudan a que la economía sea más productiva. Por lo tanto, en 2022 la economía española no cumplía esta regla. aunque se es pera que esté cerca de cumplirla en 2023 y la cumpla a partir de 2024.


\subsection{Las operaciones financieras d e las AAPP}
Las Cuentas Financieras de la Economía Española registran las operaciones y posiciones fi nancieras entre los sectores y subsectores institucionales residentes en la economía nacional y entre estos y el resto del mundo. Las operaciones financieras en el cas o de las AAPP consisten en las variaciones de activos y pasivos fi nancieros. Estas operaciones, añadidas a las no financieras , que son la parte fundamental del tema, determinan la necesidad de endeudamiento de las AAPP.

El ajuste dé.ficit-deudu\footnote{%
    Ver Forte-Carnpos el al. (2022). para los detalles, con datos de 2021 . Para los datos de 2022 y siguientes, consultar las Notas Económicas que se publiquen anualmente; así como Banco de España (2023). Ver, también, Programa de Estabilidad 2023-2026. Cuadro 3.20, página 92.
}

Si a la Necesidad (+) o Capacidad (-) de Financiación de las AAPP\footnote{%
    Cambiamos el signo habitual de esta variable y la del saldo a financiar para que los resultados de las sumas de las operacione.s aritméticas que siguen tengan signo positivo. Es lo que se hace en los Cuadros 1 1 .8, 1 1 .9 y 1 1. 1 O del Boletín Estadístico mensual del Banco de España.
} se le suma la Variación Neta de Activos Financieros, se obtiene el Saldo Neto a Financiar (+) o a Emplear (-).

Este saldo neto a financiar o a emplear equivale a la variación neta de pasivos financieros ( por ejemplo. s i el saldo es a financiar (+), la variación neta de pasivos será positiva). A su vez. la variación neta de pasivos financieros es la variación de la deuda pública PDE más unos aj ustes.\footnote{%
    Deuda pública calculada según la metodología del Procedimiento de Déficit Excesivo ( PDE) de la UE. Este es el concepto de endeudamiento relevante a efectos de los límites establecidos en el Pacto de Estabilidad y Crecimiento europeo y en la Ley de Estabilidad Presupuestaria y Sostenibilidad Financiera española. La definición actual de este endeudamiento está regulada por el Reglamento (UE) 220/20 1 4 de la Comisión. de 7 de marzo de 20 1 4
}

Nec. (+) o Cap. (-) de Fin. + Var. Neta ele AAFF = Saldo Neto a Fin. (+) o a Empl. (-) Saldo Neto a Fin. (+) o a Emp/. (-) = Var. Neta de PPFF = ti Deuda PDE + Ajustes

Las clases de activos financieros son:

1. Efectivo y depósitos. Tal es el caso de la cuenta del Tesoro en el Banco de Espai'la
2. Valores representativos de deuda.
3. Préstamos
4. Participaciones en el capital y en los fondos de inversión
5. Derivados financieros y otras cuentas pendientes de cobro

Todos ellos registrados a sus valores nominales y en términos consolidados: es decir. excluidos los activos que son pasivos de otras Administraciones Públicas.

Por su parte, los ajustes tienen diversa naturaleza; pueden ser de signo positivo o negativo. Los más notorios son los créditos comerciales y otras cuentas pendientes de pago. las emisiones por encima o por debajo del valor nominal y otros ajustes de valoración. Dependiendo de su signo. sirven para financiar el déficit (como en el ejemplo que sigue), pero no se incluyen en la deuda PDE.

En 2022 , en términos de PIB. estas magnitudes fueron:
( 1) Necesidad de Financiación de las AAPP: 4,81
(2) Adquisiciones Netas de AAFF: 3.48
(3 )=( 1 )+(2) Saldo Neto a Financiar: 8,29
(4)=(3 ) Variación Neta de Pasivos Financieros: 8,29
(4) Aj ustes: 2,63\footnote{%
    De la resta: (2) - (4): 3.48 - 2.63 = 0,85. surge el 0,9 de la fila 6 del Cuadro 3.20 (página 92), "ajuste stock-flujo''. del Programa de Estabilidad 2023-2026.
}
(5)=(3)-(4) Variación de la Deuda PDE: 5.66

Por lo tanto, la suma de operaciones financieras y no financieras I levaron a una emisión neta de deuda equivalente al 5,66 por ciento del PJB, con un stock de deuda PDE al final de 2022 del 113,2 por ciento del PIB. Cabe reseñar que en 2021 ocurrió lo contrario: la necesidad de financiación fue del 6.87 por ciento del PIB. mientras que las adquisiciones netas de AAFF menos los ajustes ascendieron al -0, 12 por ciento del PIB. Por lo tanto, la variación de la deuda PDE fue de 6,75 puntos del PIB, inferior a la necesidad de financiación.

\begin{specialblock}{\textbf{Operativa del MRR (y del PRTR)}: Contabilidad Nacional}
    El hecho de que la operativa del MRR (y del PRTR) com porte anticipos de fondos a España por parte de la Comisión Europea antes de que éstos sean gastados por las AAPP, y que en ningún mom ento se vea afectada la capacidad o necesidad de financiación de las AAPP (por el principio de neutralidad financiera), obliga a contabilizar los flujos de una manera que conviene conocer.

    Cuando la Comisión Europea realiza un anticipo a España con cargo a estos fondos, e n las cuentas financieras de las AAPP se anota una variación (positiva) de activos financieros, concretamente de la cuenta del Tesoro (o de la Administrac ión que sea) en el Banco de España, con la contrapartida en la rúbrica ·'Otras cuentas pendientes de pago'" (al Resto del Mundo).
    
    Cuando. de e��ta cantidad, las AAPP realizan un gasto (por ejemplo, de Formación Bruta de Capital (FBC)). en las cuentas financieras de las AAPP se anota una variación ( negativa) de activos financieros (con la correspondientt disminución en ·'Otras c uentas pendientes de pago'"); y en la contabilidad nacional un aumento del gasto en FBC pública y un aumento igual de ingresos de Transferencias de Capital. Simultáneamente. se anota en la B alanza de Pagos un ingreso por Transferencias de Capital de la UE y una variación (positiva) de activos financieros frente al exterior
\end{specialblock}

Después de repasar las grandes c ifras de las AAPP. nos referiremos a continuación a las principales rúbricas de gastos e i ngresos.


\subsection{Rúbricas de gastos}

IJ.lll.l. Tra11sfere11cias sociales
La rúbrica de gasto más impo11ante es la de transforencias sociales. con el 20, 1 por ciento del PIB. Pres'faciones sociales distintas de las de especie

De ellas. las más importantes son las prestaciones sociales distintas de las de especie ( 1 7,2 por ciento del PlB). Las integran, sobre todo. las pensiones de la Seguridad Social. las de Clases Pasivas del Estado, las prestaciones de desempleo y los pagos por EREs y ERTEs del SEPE , el Ingreso Mínimo Vital etc. De este porcentaje, 1 2,9 puntos son de pensiones contributivas y no contributivas de la Seguridad Social y 2.0 puntos son de las pensiones de Clases Pasivas del Estado. Los gastos de desempleo (no relacionados con la Covi d- 1 9 ) añaden 1 , 6 puntos16. Por último. destaca la Incapacidad Laboral Transitoria, con 1 ,0 puntos. Esta rúbrica es la principal responsable del incremento del gasto corriente al que nos referimos antes, por el aumento del gasto en pensiones (principalmente de jubilación) por el efecto demográfico, por la revalorizac ión de las pensiones (en especial por los aumentos que a lo largo del tiempo se han decretado para las pensiones mínimas) y por los aumentos de las bases de cotización (en línea con los salarios nominales). Entre 2008 y 2022 estos gastos aumentaron un 67,9 por ciento, m i entras que, por ejemplo, la remuneración de asalariados lo hizo en un 30,3 por ciento. Tras la reforma de las pensiones legislada a lo largo del periodo 2021-2023. que liga la revalorización de las pensiones a la inflación, ) la progresiva jubilación de la generación del ·'babyboom··. esta tendencia se acentuará. 


Tran.��ferencias .'iOCiales en especie.

De menor importancia cuantitativa son las transferencias sociales en especie. Están integradas. sobre todo. por el coste de la educación concertada y de los convenios con la sani<lad privada para la prestación de servicios. por el gasto en fannacia no hospitalana. por el pago de las mutuas de funcionarios (MUFACE, M UGEJU e ISFAS) a las sociedades de seguros de salud y por otros ser\'icios sociales (tratamientos especiales. como drogodependencia y otros que se prestan en centros privados) y determinados bienes (sillas de inválidos. etc.). De 2008 a 2022 esos gastos aumentaron un 31,9 por ciento.

ll.llI.11. Rem1111eració11 de a.\·alariado:,

La remuneración de asalariados comprende toda la remuneración en efectivo y en especie a pagar por las Administraciones Públicas a sus asalariados como conu·apartida al trabajo realizado por éstos durante el ��jercicio. Entre los asalariados se incluyen al personal de la Administración Civil del Estado. de la sanidad pública. de la educación pública. de la administración de Justicia y de las Fuerzas Annadas y de Seguridad.

Es la segunda rúbrica en impo11ancia, con el 1 1.6 por ciento del PIB. En principio, no hay indexación automática al PIB nominal o al IPC. pero. periódicamente. el Gobierno (y las autoridades regionales y locales) finnan con los sindicatos acuerdos de subidas salariales ligados al IPC y. a veces, al PIB. De ahí que no se puedan esperar caídas significativas en esta ratio.

IJ. IIJ.JJJ. Consumos intermedios

Los consumos intem1edios representan el valor de los bienes y servicios consumidos durante el proceso de producción, excluidos los activos fijos, cuyo consumo se registra en la rúbrica de consumo de capital fijo. En esta rúbrica. además de los gastos de funcionamiento de la Administración. se incluye la provisión de medicamentos en hospitales públicos (anestesias. quimioterapias. fórmulas magistrales cte.). los cuales son. generalmente, mucho más caros que los dispensados en las fannacias.

Son la tercera rúbrica en importancia. con el 5.9 por ciento del PIB. Es un gasto ligado a la intlación y a la provisión de nuevos bienes y servicios (como los nuevos tratamientos médicos que se incluyen en la cartera básica del Sistema Nacional de Salud), de ahí que no se esperen ahorros sustanciales en el futuro.

l/. III. I V. Intereses 
En esta rúbrica se recoge la carga financiera derivada de los distintos instrumentos empleados por las Administraciones Públicas en la captación de los recursos financieros necesarios para el ej ercicio de su actividad.

El gasto por intereses de la deuda ascendió en 2022 al 2,4 por ciento del PlB. Es una ratio de comportamiento anti-cíclico: baj a en las expansiones (porque al bajar el déficit baja el crecimiento de la deuda; y por el propio aumento del PJB nominal) y sube en las recesiones (porque el mayor déficit supone emitir más deuda; y por la propia disminución del PIB nominal). 

El comportamiento pro-cíclico de los tipos de interés no suele compensar esto.

IJ.III. V. Subvenciones

Las subvenciones son pagos con-ientes sm contrapartida que realizan las Administraciones Públicas. Por ejemplo, son subvenciones las ayudas y bonificaciones que se legislaron para el precio de venta al público de ciertos productos energéticos y las ayudas al sector de transporte, a causa de la guerra de Ucrania. En 2022 ascendieron al 2 por ciento del PIB. No sigue un patrón claro. pero esta ratio no tiene tendencia a subir.

Il.III. VI . Otros gastos corrien tes

Esta rúbrica contiene la apo11ación de España al Presupuesto de la UE y a los de todos los organismos de los que es miembro. así como la parte de la ayuda al desarrollo que no es transferencia de capital. Estos gastos ascendieron al 1,8 por ciento del PIB en 2022 , y no tienen tendencia.

II.JJJ . VII . Formación Bruta de Capital

Este concepto. también llamado •'inversión pública'·, comprende las adquisiciones menos las cesiones de activos fij os realizadas por las AAPP durante un período de tiempo determinado. Los activos fijos son activos producidos utilizados en la producción durante más ele un año. Es el caso de la inversión en la construcción de infraestructuras de transpo1tes (puertos. aeropuertos, carreteras, ferrocarriles) o urbanas, realizada por los ministerios. consejerías de gobiernos autónomos, por los ayuntamientos y por las sociedades que dependen de estas instituciones. También se incluyen las empresas públicas que cumplan los requisitos para ser incluidas en el sector AAPP.

Las nomrns de la contabilidad nacional obligan a ser cuidadosos a la hora de analizar el alcance económico de este agregado ya que, por ej emplo, la finalización de un periodo de concesión de una infraestr uctura y su posterior mantenimiento por las respectivas administraciones obliga a computar esta reversión como inversión pública. Y, para compensar esta operación por el lado de los ingresos ( de manera que no tenga efecto en el déficit público). se anota un incremento equivalente de las transferencias de capital. Así mismo, hay que subrayar que la adquisición de armamento militar también forma parte de e!:>ta rúbrica. así como la adquisición de material inventariable por parte de las AAPP (equipos informáticos. equipam iento hospitalario. parq ue móvil etc.)

Este agregado se situó en el 2.8 ror ciento del Pl13 en 2022. Se trata de un gasto de tipo discrecional y, por tanto. suele sufrir recortes cuando el raís está obligado a realizar aj ustes presupuestarios de fonna urgente. l��stc sería el caso de Espai'ia. Efectivamente. 2.8 puntos de PIB no está lejos del nivel de la eurozona (3.0 puntos). aunque sí está significativamente por debajo de la UE y de Estados Unidos (3,2 puntos). y de Japón ( 4.2 puntos). Históricam ente. la inversión pública en España fue muy alta a mediados de la década de 1990. alcanzando algunos años el 5 por ciento del PIB. cuando el país, con la ayuda de los Fondos Estructurales y de Cohesi ón. intentaba reducir su déficit de infraestructuras. Posteriormente. esta ratio fue disminuyendo: 4.4 puntos en el peri odo 2004-2008, 3,8 puntos en 2009-20 13 y 2.2 puntos en 2014-2018.

Tl.lll. VI II. Tr ansfe ren cias de capital
Esta rúbrica contiene dos partes: Ayudas a la inversión y Otras transferencias de capital.

Ayu das a la in versión

Las ayudas a la inversión incluyen las transferencias efectuadas por las Administraciones Públicas. en efectivo o en especie, a otras unidades instituci onales. para financ iar total o parcialmente los costes de adquisición de sus activos fijos. Son ayudas a la inversión los incentivos económ icos regionales para que las empresas invi ertan en determ inadas zonas. También lo son las transferencias a entes inversores como ADIF. ENAIRE. Pue11os del Estado. empresas productoras de amrnmento militar etc.

Otr as transferen cias de cap ital

Por su parte, las otras transferenc ias de capital recogen las transferencias que se efectúan. no para redistribuir por sí misma!> la renta, pero que sí compo11an una redistribución del ahorro o del patrimonio entre los distintos sectores de la economía. En esta parte de las transferencias de capital se encuadran las ayudas al sector financiero, las cuales tuvieron gran importancia durante la crisis financiera; por ejemplo, las transferencias del FROB a la SAREB y las indemnizaci ones a inversores minoristas. También se encuadra parte de la ayuda al desarrollo y la ayuda con armamento militar. Así mismo. cuando un préstamo de las AAPP se declara incobrable o si se ejecuta un aval público (por ej emplo, del ICO), esto se anota como una transferencia de capital.\footnote{%
    De hecho, a los de impuestos y cotizaciones sociales devengados se les aplica anualmente un ajuste por recaudación incierta. Este ajuste incrementa esta rúbrica de gasto.
} También la ejecución de senten cias judiciales y laudos imponi endo sanciones y obligando a la Administración a devolver impuestos, tasas o cánones indebidamente recaudados, se anotan en esta partida.

Esta rúbrica ascendió en 2022 al 1 . 1 por ciento del PIB y no tiene tendencia a crecer.

\begin{specialblock}{\textbf{Gasto público computable:} Regla del gasto}
    La legislación macro-fiscal de la UE, que entrará en vigor en 2024, tendr,i como variable objetivo el nivel de deuda pública. como porcentaje del PJB; y como variable instrumental, una medida del gasto público, como porcent��je del PIB. El stock de deuda pública y su trayectoria es la magnitud más relacionada con la sostenibilidad fiscal. Y el nivel de gasto público .y su tasa de crecimiento es el principal determinante de la deuda y de su trayectoria a medio y largo plazo. El cumplimiento o 110 de la regla de gasto, en las distintas variantes que veremos, se verifica con los datos de i ngresos y gastos de la ��uenta de las AAPP que estamos explicando en este tema.\footnote{%
        Los gastos de desempleo no figuran desglosados en esta clasificación de gastos. Es preciso recurrir a la clasificación COFOG que veremos a continuación.
    }
    
    En la UE, el gasto público objeto de control es aquél que no está financiado con fondos europeos, puesto que el financiado con fondos de la UE no influye en el déficit público.\footnote{
    
    } A su vez, a esta magnitud se le aplican las siguientes correcciones antes de obtener la medida de gasto a controlar mediante la regla fiscal.
    
    Una primera medida del gasto público que se utiliza en este tipo de reglas es el [!asto púhlico primario neto de medidas discrecio11ales de inRresos; de manera que, en ausencia de obligaciones de ajustes adicionales, el gasto público primario no podrá aumentar más que el PlB potencial nominal, añadiendo el aumento de recaudación derivado de medidas de ingresos que hayan sido legisladas. Estas magnitudes de gasto y de ingreso son las que figuran en una cuenta como la de la tabla 1. es decii:, en términos de contabilidad nacional. 

    Otro retoque adicional a esta magnitud lleva a utilizar el gasto primario neto de medidas discrecionales de ingresos y neto de los gastos cíclicos de desempleo. Así, a la med ida inicial del gasto primario se le resta cada año el gasto de desempleo que excede del gasto correspondiente a la tasa estructural de desempleo (la NA WRU). Esto, cuando el output gap es negativo. Cuando es positivo, hay que sumar al gasto primario los gastos de desempleo que no han ocurrido por ser el PJB observado mayor que el PIB potencial ( y la NA WRU mayor que la tasa de paro observada U'). Esto se puede calcular de varias fom1as. Últimamente, en la U E se calcula la desviación de la NAWRU con respecto a la tasa de desempleo observada, y se multiplica por el gasl0 total de desempleo (GTD): GCD = u·-N u��WRU x GTD. La cantidad obtenida ( GCD) es el gasto cíclico de desempleo. Si es positivo. se resta del gasto primaiio. Si es negativo, se suma al gasto primario.

    En la regla fiscal, la variable de gasto público primario elegida se puede retocar todavía más. de manera que no se socaven otros objetivos de las políticas públicas. Por ejemplo:
    1. No computar como incremento de gasto detcnninadas inversiones públicas
    2. No computar como incremento de gasto las derivadas de medidas para implementar la transición verde
    3. No computar como incremento de gasto el gasto militar adicional derivado de los compromisos con los países de la OTAN y sus aliados

    En el caso de un país altamente descentralizado, como Espai'ía. el cumplimiento de la regla fiscal depende del desempefío de los subscctores de las AA PP. Por ello, se suelen aprobar reglas fiscales a nivel sub-nacional para garantizar el cumplimiento a nivel nacional. La verificación del cumplimiento o no de las reglas por parte de los sub- sectores de las AAPP se verifica, igualmeme. con sus respectivas cuentas, que veremos más adelante.
\end{specialblock} 


\subsection{Rúbricas: Ingresos}
Sabemos que los ingresos provienen principalmente de impuestos como el lRPF. el Impuesto sobre Sociedades, el IV A y los Impuestos Especiales. Todos ellos son estudiados en otra parte del temario. Sin embargo, hay otros muchos impuestos, tanto estatales como autonómicos y locales que no sólo tienen impo11ancia c uantitativa. sino que son esenciales para lograr el objetivo de saldo público. De hecho, algunos se han creado para afiadir más recursos a los ya apo11ados por los grandes impuestos. Así mismo. hay ingresos no tributarios que es preciso tener en cuenta. A diferencia de los gastos. que en ausencia de indexación son relativamente rígidos (a excepción de los gastos de desempleo, que no tienen tendencia y son contra-cíclicos), los ingresos, a medio y largo plazo, suelen evolucionar en paralelo con el PIB nominal, si bien las subidas y bajadas discrecionales de impuestos pueden introducir irregularidades en los datos.

II.IV.J. Impuestos indirectos

Esta rúbrica incluye aquellos impuestos que gravan la producción y la importación de bienes y servicios, la utilización de mano de obra. la propiedad o el uso de la tierra, los edificios y otros activos utilizados en la producción.

Esta rúbrica se divide en dos conceptos: lmpuestos sobre los productos y Otros impuestos sobre la producción. La lista de impuestos indirectos es muy larga. A continuación, se relacionan los más importantes: 

Impuestos sobre los productos
Se definen como aquellos impuestos a pagar por cada unidad producida o distribuida de un bien o servicio.

Por una parte. están los impuestos tipo IVA, como el IV A y el Impuesto General Indirecto Canario ( IGJC). y los Derechos a la impo1tació11,

Por otra parte, los impuestos sobre los productos, excluidos el IV A y los impuestos sobre las importaciones, comprenden los impuestos sobre los bienes y servicios que gravan la producción, exportación, venta, transferencia, anendamiento o entrega de dichos bienes y servicios o su utilización para consumo final propio o para la formación de capital por cuenta propia. En este concepto se encuadran: los Impuestos Especiales, el Impuesto sobre Transacciones Financieras, el Impuesto sobre las Primas de Seguro, el Impuesto Especial sobre Determinados Medios de Transporte (conocido como ''impuesto de circulación''), la Tasa sobre el Juego. el Impuesto sobre Detem1inados Servicios Digitales y el Impuesto sobre Construcciones, Instalaciones y Obras ( ICIO).

En 2022. los ingresos de esta rúbrica fueron los terceros en importancia. después de los impuestos directos y las cotizaciones sociales, con el 12, 1 por ciento del PIB. Los impuestos tipo IV A fueron el 7.1 por ciento del PIB y los impuestos excluidos el IVA y los derechos de importación fueron el 3,2 por ciento del PIB.

Otro ... impuestos sobre la producción

Recogen la imposición soportada por las empresas como resultado de su participación en la producción, independientemente de la cantidad o el valor de los bienes y servicios producidos y vendidos. Incluyen: el Impuesto sobre Bienes Inmuebles (IBI),\footnote{%
    Efectivamente: a pesar de que el SEC-201 O dice que se trata de "impuestos que soportan las empresas", hay en esta rúbrica impuestos que pagan los hogares. como el .1 BJ. las licencias urbanísticas o el canon ITV.
} el Impuesto sobre Vehículos de Tracción Mecánica (cuando recae sobre las empresas). el Impuesto sobre Actividades Económicas (IAE). las Licencias urbanísticas, el Canon ITY, los Derechos de Emisión de Gases de Efecto Invernadero y el recurso ordinario del Fondo de Garantía de Depósitos. En 2022, estos impuestos recaudaron el equivalente al 1,8 por ciento del PlB.


/1.JV.IJ. Impuestos directos
Esta rúbrica comprende todos los pagos obligatorios sin contrapartida, en efectivo o en especie. recaudados periódicamente por las Administraciones Públicas sobre la renta y el patrimonio de las unidades institucionales. así como otros impuestos periódicos corrientes.

Los principales impuestos son: el IRPF, el Impuesto sobre Socied ades, el Impuesto sobre la Renta de No Residentes. el Impuesto sobre el Patrimonio, el IB1 (de viviendas desocupadas) y el Impuesto sobre Vehículos de Tracción Mecánica (cuando recae sobre los hogares).

Estos impuestos fu eron en 2022 los segundos en importancia después de las cotizaciones sociales. con el 12,4 por ciento del PIB. De largo. los más importantes fue ron los impuestos sobre la renta. con el 12. l por ciento del PIB. Dentro de los impuestos sobre la renta. la distribución es, aproximadamente, 80 por ciento sobre la renta de personas físicas y 20 por ciento sobre la renta de sociedades (residentes y no residentes, en ambos casos).

Coti:acio11es sociales

Esta rúbrica recoge las cotizaciones sociales efectivas o imputadas que pagan los hogares a los sistemas de seguros sociales con el fin de asegurar el pago de prestaciones sociales. Las cotizaciones sociales efectivas son en su mayor parte los pagos obligatorios que los empleadores. asalariados. trabajadores por cuenta propia y desempleados realizan a los Fondos de la Seguridad Social, empresas de seguro y otros sistemas de seguros sociales, con el fin de garantizar el <lcrecho a las prestaciones sociales a sus asalariados o bien, en su propio beneficio, como en el caso de los trabaj adores por cuenta propia y los desempleados.

Las cotizaciones sociales imputadas representan la contrapartida de las prestaciones sociales pagadas directamente por los empleadores a sus asalariados. antigu os asalariados y otros derechohabientes. Por ej emplo, en el Régim en de Clases Pasivas del Estado (en extinción). durante un determinado afio, las cotizaciones sociales imputadas equivalen a la diferencia entre las pensiones pagadas a los jubilados y las cotizaciones efectivamente ingresadas, pagadas por los funcionarios en activo. La palabra "imputadas'· quiere decir que es el Estado. como empleador. el que en la pr áctica aporta las cotizaciones necesarias para completar lo pagado por los funcionarios y, así. poder pagar a los jubilados.

En 2022 fue la primera rúbrica en importancia. con el 13.6 por ciento del PlB. Las cotizaciones sociales efec tivas a cargo de los empleadores ascendieron al 9,6 por ciento del PIB. Las cotizaciones sociales efectivas a cargo de los hogares ascendieron al 3,4 por ciento del PIB. Y las cotizaciones sociales imputadas ascendieron al 0,5 por ciento del PIB.

ll. JV.Jll. Otros ingresos corrientes
Esta rúbrica contiene ingresos de diversos orígenes. Fundam entalmente, son rentas de la propiedad, ingresos por ventas de bienes y servicios, y transferencias coITientes. 

Las rentas de la propiedad recogen los recursos recibidos por las Administraciones Públicas como propietarias de activos financieros, a cambio de proporcionar fondos a otras unidades institucionales. Se trata. por ej emplo. de los beneficios del Banco de Espafia transferidos al Estado. de los dividendos de las em presas públicas. de los ingresos por venta de loterías. etc. 

Los recursos por la venta de bienes y servicios. pueden provenir de la venta de producción de mercado o la venta de producción de no mercado. En el primer caso. es aquella que ��e vende en el mercado a precios económicamente significativos o se conserva para su consumo final propio o para la formación de capital. En el segundo caso, es aquella producción suministrada a otras unidade s de forma gratuita o a precios económicamente no s i gnificativos : proceden en su mayor parte de los pagos relacionados con la educación, los servicios sociales, el deporte y las entradas a museos.

Por último, las transferencias corrientes incluyen los recursos de la cooperación internacional corriente, las ifüfomnizaciones de seguro no vida y las transferencias corrientes diversas. Las más i mportantes son las primeras. Comprenden todas las tran sferencias, en efectivo o en especie, entre las Administraciones Públicas del Resto del Mundo y las Administraciones Públicas nacionales. con excepción de las ayudas a la inversión y las otras transferencias de capital. Se trata, por ejemplo, de las transferencias provenientes del Fondo Social Europeo (FSE). Por su parte, las transferencias corrientes diversas proceden, básicamente. de reintegros por operaciones corrientes y de capital, de recargos, sanci ones y multas tributarios -siempre que el i mporte de estos conceptos tributatio s pueda diferenciarse de los propios i mpue stos- y no tributario s, y de recursos eventuales.

Como puede verse en la tabla l , estos recursos ascendieron en 2022 al 3.8 por c iento del PlB. aproxi madamente lo mismo que en años anteriores.

JI.IV.IV. Impuestos sobre el capital

Los i mpuestos sobre el capital gravan, a intervalos irregulares el valor de los activos o e l patrimonio neto d e las unidades institucionales o e l valor d e los activos transferidos entre unidade s institucionales , como resultado de sucesiones, donaciones inter-vivos u otras transferencias. Estos impuestos que recogen, entre otros. los recursos derivados de los gravámenes sobre sucesiones y donaciones, e l impuesto sobre el incremento del valor de los terrenos de naturaleza urbana. las cuotas de urbanizaci ón, los aprovechamientos urbanísticos, las contribuciones especiales y el recurso extraordinario del Fondo de Garantía de Depósitos: ascendieron en 2022 al 0,4 por ciento del PIB. que viene a s er su nivel nonnal.



\clearpage
\section{Estructura de gastos: Clasificación funcional del Gasto (COFOG)}
\puntosclave{
    La clasificación COfOG divide el gasto público en 1 O funciones. 
    
    • El gasto social. imegrado. principalmente, por las funciones de Protección Social. Sanidad y F.ducación es el más importante (69.6 por ciento del total), seguido de Servicios públicos generale!> y de Asuntos Cconómicos. 
    
    La carga o presión fiscal global de un país son los impuestos y cotizaciones sociales recaudados por las administraciones públicas y las instituciones europeas. así como los impuestos que pagan los sectores residentes al resto del mundo. 
    
    En 2022 la presión fiscal en España era del 38,2 por ciento del PIB. mientras que en la Unión Europea era del 40.4. En general, Espai'ia aparece entre los países desarrollados como un país con un nivel de presión iiscal medio. 
    
    J\l igual que en países como /\lemania y Francia, y a diferencia de los países anglosajones. en España tienen gran peso las cotizaciones sociales. 
    
    Dentro de las cotizaciones sociales, tienen mayor peso aquellas a cargo de los empleadores (72 por ciento del total). En el caso de los impuestos sobre la renta y los beneficios, predominan los que recaen en la renta de las personas físicas. con el 81.6 por ciento. 
    
    La Administración Central y los FFSS son los que generan más presión fiscal. 
    
    Por función económica. destacan los impuestos sobre el trabajo (50.8 por ciento del total)
}

\begin{specialblock}{\textbf{TODO ESTO ESTABA EN EL ANAXO}}
    La Clasificación de Funciones de las Administraciones Públicas (COFOG) es una estadística desarrollada por la OCDE y publicada por la División de Estadísticas de las Naciones Unida (UNSD) que clasifica el gasto, calculado de acuerdo con la metodología del SEC-20 l O, conforme a la fi nalidad que persiguen los fondos,
    
    La COFOG estructura el gasto en varios niveles: un primer nivel, con l O divisiones de gasto (grupos funcionales); un segundo nivel, en el que se desaITolla cada una de las divisiones anteriores en un máximo de 9 grupos (sub- funciones); y un tercer nivel, en el que se desagrega en clases cada uno de los grupos anteriores, aunque este nivel de desagregación no está implementado en todos los países, como ocurre en el caso de Espai'ía.

    01 Servicios públicos generales
    O l. l Órganos ej ecutivos y legislativos, asuntos financieros y fiscales. asuntos exteriores
    01.2 Ayuda económica exterior
    O 1.3 Servicios generales
    01 ,4 Investigación básica
    01.5 l+D relacionada con los servicios públicos generales   
    01.6 Servicios públicos generales n.c.o,p,
    01.7 Operaciones de deuda pública
    01.8 Transferencias de cankter general entre distintos niveles de las admones. públicas
    02 Defensa
    02 . 1 Defensa militar
    02.2 Defensa civil
    02 .3 Ayuda militar al exterior_
    02.4 I+D relacionada con la defensa
    02 ,5 Defensa n.c,o,p,
    03 Orden público y segurida d
    03.1 Servicios de policía
    03.2 Servicios de protección contra incendios
    03,3 Tribunales de justicia
    03.4 Prisiones
    03,5 l+D relacionada con el orden público y la seguridad
    03,6 Orden público y seguridad n.c.o,p.
    04 Asuntos económicos
    04 .1 Asuntos económicos. comerciales y laborales en general
    04.2 Agricultura, silvicultura, pesca y caza
    04,3 Combustible y energía
    04.4 Minería, manufacturas y construcción
    04 ,5 Transporte
    04,6 Comunicaciones
    04. 7 Otras acti\ idades
    04.8 l+D relacionada con asunto!> económicos
    04.9 Asuntos económicos n.c.o.p.
    05 Protección del medio ambicntr
    05.1 Gestión de los residuos
    0 5.2 Gestión de aguas residuales
    0 5.3 Reducción de la con1aminación
    0 5.4 Protección de la biodiversidad y del paisaje
    05.5 l+D relacionada con la protección del medio ambienle
    05.6 Protección del mcdioambicnte n:c.o.p.
    06 ViYienda y servicios comunitarios
    06. 1 Urbanismo
    06.2 Desan-ollo comunitario
    06.3 Abastecimiento de agua
    06.4 Alumbrado público
    06.5 l+D relacionada con la vivienda y los sen·icios comunitarios
    06.6 Vivienda y servicios comunitarios n.c.o.p.
    07 Salud
    0 7.1 Medicamentos y otros productos farmacéuticos. aparatos y material terapéutico
    0 7.2 Servicios ambulatorios
    0 7.3 Servicios hospitalarios
    0 7.4 Servicios de salud r,ública
    07.5 I+D relacionada con la salud
    0 7.6 Salud n.c.o.p.
    08 Ocio, cultura y religión
    08.1 Servicios recreativos y deportivos
    08.2 Servicios cul!urales
    08.3 Servicios de radio y televisión y servicios editoriales
    08.4 Servicios religiosos y otros servicios comunitarios
    08.5 I+D relacionada con el esparcimiento, la cultura y la rcl igión
    08.6 Actividades recreativas, cultura y religión n.c.o.p.
    09 Educación
    09.1 Educación preescolar y primaria
    09.2 Educación secundaria
    09.3 Educación postsecundaria no terciaria
    09.4 Educación terciaria
    09.5 Educación no atribuible a ningún nivel
    09.6 Servicios auxiliares de la educación
    09. 7 I+D relacionada con la educación
    09.8 Educación n.c.o.p.
    10 Protección social
    10.1 Enfermedad e incapacidad
    10.2 Edad avanzada
    10.3 Supérstites
    10.4 Familia y hijos
    10.5 Desempleo
    10.6 Vivienda
    1 O. 7 Exclusión social n.c.o.p.
    10.8 l+D relacionada con la protección social
    1 0.9 Protección social n.c.o.p.
    n.c.op.: No clasificada en otras partidas
\end{specialblock}


La contabilidad nacional agrupa el gasto de las AAPP en las 1 O funciones de la clasificación COFOG (Classification of the Functions of Government) que aparecen en la tabla 2 :

\textbf{TABLA RESUMEN: DISTRIBUCIÓN COFOG DEL GASTO (\%)}

El gasto social. englobado en las funciones de Protección Social. Sanidad, Educación, Actividades recreativas, cultura y re ligión; Medio ambiente y Vivienda y (por este orden). es. con el 33 por ciento del PIB e l área más destacada del gasto (69,6 por ciento del total). Destaca. también. la relativa estabilidad de los porcenta,jes de gasto de cada función.

En IGAE (2022a) hay información desglosada por niveles de gobierno. Se opta por no incluirla en el tema por falta de tiempo.

\subsection{La presión fiscal}

La carga o presión fiscal global de un país son los impuestos y cotizaciones sociales recaudados por las administraciones públicas y las instituciones europeas. así como los impuestos que pagan los sectores residentes al resto del mundo. En general. la presión fiscal se expresa en función del PI B, es decir, se m ide en términos de la relación entre impuestos y cotizaciones sociales y e l PIB.\footnote{%
    El PIB se suele r.::visar a lo largo de los años. Por ello, la ratio de la presión fiscal pude cambiar retroactivamente cuando esto sucede.
}

Los impuestos y cotizaciones sociales que han de registrarse para calcular la presión fiscal son los impuestos sobre la producción y las importaciones (impuestos indirectos). los impuestos corrientes sobre la renta, el patrimonio ( impuestos directos), etc., las cotizaciones sociales netas y los i mpuestos sobre el capital, tal y como se definen en e l SEC 2010 y con los criterios de registro en él establecidos.\footnote{%
    Esto es importante. Con relativa frecuencia aparece en la prensa información sobre presión fiscal que no coincide (a veces por bastante d iferencia) con los datos de organismos a los que nos referiremos a continuación. Esto ocurre si se toman magnitudes de recaudación en vez de contabilidad nacional. También es i mportante respetar la metodología correcta.
} Se incluyen los impuestos recaudados por la UE y. en su caso, el ajuste por recaudación incierta de impuestos y cotizaciones sociales.

Es una magnitud relevante para realizar comparaciones internacionales. Tanto la Comisión Europea2.\footnote{%
    La metodología es la del Reglamento Delegado U E nº 877/20 1 3 de la Comisión. de 27 de junio de 2 0 1 3. por el que se completa el Reglamento (UE) 11º 473i20 1 3 del Parlamento Europeo y del Cons��jo sobre disposiciones comunes para el seguimiento y la evaluación de los proyectos de planes presupuestarios y para la corrección del déficit excesivo de los Estados m i embros de la zona euro.

Además, como se verá, dependiendo de los datos que sean, el último publ icado puede ser de 2022. 202 1 o 2020.
} ( ver Comisión Europea ( 2023a) y ( 2023b)). como la OCDE ( ver OCDE (2022) ), ¡rnblican datos y estudios sobre presión fiscal. Utilizaremos una u otra fuente en función de la información que proporcionen25. Las metodologías no son plenamellle coincidentes. pero los resultados son muy parecidos.

JJJ.Il.l. La presión jisml en Espmia: comparación internacional

En los datos de la base AMECO. en 2022 España aparece con una presión fiscal del 38.23 por ciento del PIB.\footnote{%
    En la Actualización del Programa de Estabilidad 2023-2026 (publicado un mes antes), se consignaba la cifra de 38,7 por ciento del PJB.
} La Unión Europea tenía una presión fiscal del 40.42 por ciento y la U E - 1 5 del 40.77 por ciento. España es. pues. un país con una fiscalidad relativamente baja en Europa.

Sin embargo. tiene interés extender la comparación a otros países. Para ello tenemos que recurrir a los datos de la OCDE (2022), que llegan hasta 202 1 . En la rabia 2 se compara la presión fiscal de varios países de referencia:\footnote{%
    La OCDE elabora estos datos a partir de la información que le envían las delegaciones nacionales. Puede haber diferencias metodológicas, especialmente entre los países de la UE y los demás.
}

\textbf{TABLA RESUMEN: PRESIÓN FISCAL (OCDE), COMPARATIVA \%}

La tabla 3 está encabezada por dos países escandinavos (Dinamarca y Suecia). más Francia: conocidos por la fuerte presencia del Estado en sus economías. A continuación. Alemania y España están algo por debajo de la media de la UE. El siguiente grupo es el de los países anglosajones (Nueva Zelanda, Reino Unido, Australia y Estados Unidos). Cierran la tabla dos países hispano-americanos (Chile y México), además de Irlanda. que es un auténtico "atípico·• en materia fiscal, teniendo en cuenta su alta renta per cápita.

Volviendo a España. llama la atención que la subida de la presión fiscal con respecto a 2019 haya sido de 3,4 puntos; mientras que en la UE registró una caída de O, 12 puntos y en la UE-15 la caída fue de 0,91 puntos. La Actualización del Programa de Estabilidad 2023-2026 (Cuadro 3.11, página 74) prevé que la presión fiscal se sitúe en 2024 por encima del 40 por ciento del PIB. y que seguirá subiendo.

El aumento de la presión fiscal se ha debido a una evolución de los ingresos tributarios marcadamente mejor que la de las magnitudes nominales de la economía. Así. con datos de recaudación de la Agencia Tributaria. la demanda interior creció. en términos nominales, un -11 ,1 por ciento en 2020, un 8,7 por ciento en 2021 y un 12,0 por ciento en 2022: mientras que la recaudación total de impuestos lo hizo en un -8,8, en un 15,1 y en un 14,4 por ciento. en esos mismos ai\os. También ha habido subidas de tipos impositivos; por ejemplo, en el lRPF a las rentas altas; pero ha habido también bajadas de impuestos y aumentos de bonificaciones.

Desde luego. un aumento ele la presión fiscal de esta magnitud no se explica por una suhida de la fiscalidad a no ser que se extienda a todos los impuestos, cosa que no ha ocurrido. En los rec;uadros 1 y 2 del Plan Presupuestario 2023 (ver bibliografía) se menciona el afloramiento de bases imponibles durante la crisis de la Covid-19. quizás para poder optar a las ayudas que se pusieron en pie. También se alude a la bajada del desempleo estructural. Con estimaciones de la Comisión Europea de la primavera de 2023. la NA WRU en España se situó en el 11.0 por ciento de la población activa. mientras que en 2019 era el 12.9 por ciento.


IJl.fl.ll. La presión fiscal en Espmia: desglose por tipos de impuestos

En la tabla 4 se muestra. con datos de la OCDC (2022) la distribución de la pre��ión fiscal en Espa1ia en 2021 por tipo de impuestos. y se compara con algunos países de referencia.

\textbf{TABLA RESUMEN: PRESIÓN FISCAL POR TIPOS DE IMPUESTOS}

En el caso de España. la tabla 4 refleja la impo1tancia que tienen las cotizaciones a la Seguridad Social (35.6 por ciento del total) en la estructura fiscal; así como la importancia algo mayor de los impuestos sobre la renta y la propiedad, en relación con los impuestos indirectos. Este es el patrón que siguen otros países como Francia y Alemania.

En cambio. en Dinamarca las cotizaciones sociales no tienen ningún peso.\footnote{%
    Efectivamente. el sistema básico de pensiones públicas en Dinamarca se financia con impuestos y cubre a toda la población residente en el país. A su lado, el sistema ocupacional es obligatorio para los trabajadores por cuenta ajena, es de contribución definida y se nutre de aportaciones de empleadores y empleados que no tienen carácter de cotizaciones sociales.
} Este patrón se repite en Australia y Nueva Zelanda. Irlanda. Reino Unido, Suecia y Estados Unidos destacan también por el bajo peso relativo de las cotizaciones sociales. Todos estos países prefieren cargar el peso de la fiscalidad en los impuestos sobre la renta y la propiedad. En el caso de Estados Unidos. llama también la atención la baja fiscalidad sobre el consumo.

Esto puede ser importante para el funcionamiento de la economía. De hecho. las cotizaciones sociales son sospechosas de perjudicar el empleo. al ser un impuesto a la mano de obra. Así, Prescott (2004) afirmaba que los impuestos explican la mayor parte de la diferencia e11tre Estados Unidos y Europa en la evolución de las horas trabajadas.\footnote{%
    Prescott. E. C. (2004) ·'\Vhy Do Americans Work So Much More Than Europeans? Fed,m1/ Res.:rve Bank of.\fi1111eupolis Q11ar1er��v Revie11'. 28. 2-1 3.
}


\textbf{Presión fiscal de España en detalle (mirar apuntes): RENTA Y BENEFICIOS / COTIZACIONES / RIQUEZA / BIENES Y SERVICIOS}

\textbf{Presión fiscal de España en detalle (mirar apuntes): Por NIVELES DE GOBIERNO}

\textbf{Presión fiscal de España en detalle (mirar apuntes): Por FUENTES/USOS DE RENTA}



\begin{specialblock}{\textbf{Trabajo:} Cálculo de la cuña fiscal}
    La fiscalidad del trabajo es especialmente relevante por los efectos que puede tener sobre el mercado de trabajo y sobre el crecimiento económico (ver Comisión Europea (2023b). página 41 ).
    
    A este respecto. tiene importancia la cuña fiscal del trabajo. que surge de la diferencia en1re el salario bruto pagado por el empleador y el salario neto del trabaj ador-1 ��. Si llam amos W al salario pactado entre el empresario y el trabaj ador, el coste salarial que paga el empresario es: W (1 + es E): es decir. el salario más la cotización social a cargo del empresario. En tanto que el trabajador percibe: W(l - r1 )(1 - csy ); es decir, el salario pactado menos el lRPF sobre las rentas del trabajo (siendo r1 el tipo impositivo medio sobre las rentas del trabajo), y menos la cotización social a cargo del trabajador. Dividiendo la segunda expresión entre la primera y multiplicando por 100, obtenemos: W(l -r¡ )(l -CSy ) ( l -r¡)(l - csr ) ------- X 100 = ------ X 100 W(l + csE ) (1 + csrJ
    
    Esta expresión es la retribución neta percibida por el empleado como el porcentaje del coste salarial LOtal soportado por el empleador para el que trabaja.

    La cui'la fiscal, .O., es 100 menos esta expresión; es decir, el porcentaje del sala1io bruto pagado por el empleador que no es percibido por el trabajador a causa del pago de impuestos y cotizaciones. : (1 - r1 )(1 - csr) .0. =100 - ( ) xl00 1 + CSr-;
    
    Según la Comisión Europea (2023a), la cuña fiscal en España en 2021 fue del '27,9 por ciento para un trabaj ador soltero, sin hijos. que gana el sahu-io medio. En la lJE fue del 31 , 9 por ciento. Esto significa que el 27 .9 por ciento del coste laboral en España se debe a la fis calidad y a las cotizaciones sociales.
    
    Un refinamiento adicional de esta medida sería tener en cuenta el poder adquisitivo de lo pagado por el empresario y lo percibido por el trabajador. Se trata de introducir en la ecuación la im posición (indirecta) sobre el consumo. Efectivamente, lo que más le imeresa al trabajador en una negociación salarial es lo que puede comprar con el salario que le quede disponible para consumir. 

    Así. el salario real pagado por el empresario es:    Wr, = p
    
    Y el salario real percibido por el trabajador es: Wr = W(l - T¡ ) ( l - esT ) 
    
    Siendo Te el tipo impositivo medio sobre el consumo y P el índice de precios de producción.

    Por lo tanto, la cuña fiscal modificada sería: w ( 1 - r )(1 - es ) fl' = 100 - 2 X 1 00 = 100 - I T X 1 00 wE ( 1 + Tc)(1 + cs,J

    Por lo tanto. la imposición sobre el consumo es relevante para el mercado de trabajo. En el Modelo EREMS,\footnote{Ver Boscá et al. (20:20).
    } el valor calibrado para Te en España es 0.15. Tomando este dato, la cuJia fiscal modificada sería el 37,3 por ciento en España.\footnote{%
        Estos datos de cuña fiscal utilizan tipos impositivos -medios. Pero. como señala Prescott ( 2004). lo relevante para la oferta de trabajo son los tipos impositivos marginales. Utilizando su metodología, la cuña fiscal en España sería mucho más alta. superior al 45 por ciento en el caso de la cuña fiscal modificada (ver Boscá et al. (2022)).
    }
\end{specialblock}




\clearpage
\section{Cuentas subsectores}
\puntosclave{
    Dentro de la Necesidad de Financiación de las AAPP destaca la de la Admini!>lración Central. Esto se debe a las transferencias a las demás adm inis1raciones, garantizando la suficiencia financiera de éstas para el cumplimienlO de sus funciones. 

    • En \ Olumen de gasto tiene más importancia la Administración Central. con el 50 .2 por ciento del total. seguido de los fondos de la Seguridad Social. con el 33.8 por ciento.

    • La rúbrica de gasto más importante de la Administración Central ,;on las Transferencias en'tre AAPP (el 52 por ciento de su gasto total), por lo dicho en el primer punto.
    
    • En remuneración de asalariados ) en consumos intermedios. tiene más imrortancia la Administración Regional. Y en Transferencias sociales. los FFSS

    • El esfuerzo inversor está bastante repartido entre la Administración Central y las Territoriales.
    
    • Las cuentas financieras ponen de relieve la importancia de la AC dentro del conjunto de la Administración y de la economía. Especialmente, por endeudarse para suministrar financiación a las administraciones territoriales y a la Seguridad Social.
}

Después de repasar las operaciones no financieras de las AAPP consideradas en su conjunto. pasamos a ver cómo e tán distribuidas estas operaciones entre las distintas Administraciones (Central, Regional_ Local y Fondos de la Seguridad Social), y las interrelaciones que hay entre ellas. Los subsectores son:
• Administración Central (AC)
o Estado
o Organismos de la Administración Central

• Administración Regional (AR)

• Administración Local (AL)

• Fondos de la Seguridad Social (FFSS)\footnote{%
    En con1c1bilidad presupuestaria el SEPE y el FOGASA son Organismos Autónomos del Estado, pero en contabilidad nacional son Fondos de la Seguridad Social. El FOGASA tiene como misión principal abonar a los trabajadores los ��alarios pendientes de pago a causa de insolvencia. suspensión ele pagos, quiebra o concurso de acreedores de los empresario��. La misión principal del SEPE es abonar las prestaciones por desempleo.
}
o Sistema ele la Seguridad Social
o Servicio de Empleo Público Estatal (SEPE) 
o Fondo de Garantía Salarial (FOGASA)

En el caso de la AC, sus Organismos incluyen un gran número de unidades públicas no de mercado que adoptan diferentes formas jurídicas (Organismos Autónomos. Entidades de Derecho público, Sociedades mercantiles, Fundaciones, Consorcios. Universidad Nacional de Educación a Distancia) y que por su actividad y naturaleza económica se consideran, a efectos de la contabilidad nacional. dependientes de la Administración Central. A pesar de que la IGAE proporciona información desglosada de la AC y de sus Organismos, nos limitaremos a dar los datos consolidados, pues los Organismos son en realidad instrumentos de la AC para realizar sus funciones.

Del mismo modo, las Administraciones Regional y Local incluyen las unidades públicas no de mercado controladas por ellas. En el caso de los Fondos de la Seguridad Social (Sistema de la Seguridad Social, SEPE y FOGASA). se incluyen unidades públicas que gestionan prestaciones sociales, como las m utuas.

Una cuestión importante es el encuadramiento o no del sector público empresarial en el sector AAPP ( tanto en la Administración Central, como en la Regional o en la Local). Para dio. se sigue el criterio del SEC 20 l O, de manera que para que un ente o empresa se encuadre dentro de las AAPP debe cumplir un doble requisito: (i) que esté bajo el "control" de las administraciones públicas y (ii) que sea un productor "no de mercado''. Se entiende que existe ·'control'' si no tienen una clara independencia en el ejercicio del poder de decisión económica. Se dice, que son "no de mercado'' si proporcionan todos o pa1te de sus bienes y servicios gratis o a precios que no son '·económicamente significativos.\footnote{%
    Los criterios para determinar sí existe control, i ndependencia y el que la sociedad sea o no de mercado, son bastante prolijos y van más allá del interés de este tema. En el caso de España. en la publicación (IGAE 2022b) que figura en la bibliografía, se recoge l a lista de todas las entidades que pertenecen al sector AAPP
}

Aplicando estos criterios, las entidades encargadas de elaborar la contabilidad nacional (INE, lGAE y Banco de Espafia) determinan cada afio (IGAE (2022b) qué sociedades se encuadran en las AAPP y cuáles no. En la actualidad, quedan fuera de las AAPP grandes sociedades inversoras (Puertos del Estado, RENFE. ENAIRE, ADIF-Alta Velocidad etc.).\footnote{%
    En cambio, ADIF (Infraestructuras) sí está dentro ele las AAPP. porque al no cubrir su actividad comercial el 50 por ciento de sus costes (necesitando, por tanto. transferencias del Estado), se considera que sus precios no son económicamente significativos
}


\subsection{Los saldos presupuestarios de los su bsectores de las AAPP}

Al final de la tabla 1 aparecía la capacidad o necesidad de financiación de las AAPP. En la tabla 6 se muestra el reparto de este saldo entre las distintas administraciones

\textbf{TABLA RESUMEN: CAP/NEC AAPP ESPAÑA}

La tabla 8 muestra que en 2019 las AAPP en su conjunto estaban a punto de llegar al nivel de déficit del 3 por ciemo del PIR que en la legislación de la UE marca el nivel del déficit excesivo. Llama la atención los superávits de la Administración Local. De hecho, los ayuntamientos y diputaciones han venido registrando superávits en torno a medio punto del PIB desde 2012. Esto es así porque la legislación de estabilidad fiscal española considera que todos sus gastos e ingresos son de naturaleza estructural; además. les obliga a tener sus cuentas en equilibrio o superávit. Y. por último. están sometidos a la regla de gasto (al igual que las demás administraciones ). Esto último hace que. una vez que obtienen un superávit, al aplicar la regla de gasto el superávit se mantiene o se amplía.\footnote{%
    En respuesta a la solicitud de la Federación Nacional de Municipios y Provincias, en ciertos años las leyes de Presupuestos han relajado algo la regla de gasto, excluyendo del cómputo algunos gastos, por ejemplo, en infraestruc1uras: o permitiendo a los ayuntamientos incorporar créditos no ejecutados de años anteriores.
}

El aumento del déficit provocado por la crisis de la Covid-19 (la regla fiscal no se aplicó en el periodo 2020-2021) fue soportado sobre todo por el Estado (que está incluido en la Administración Central). que suministró financiación a las demás administraciones. como veremos más adelante.

En 2022 los desequilibrios · fiscales bajaron a unos niveles menos preocupantes. La Administración Central registró un saldo de -3.09 por ciento del PIB. Esto fue producto de una Necesidad de Financiación del Estado del 3,31 por ciento del PIB y de una Capacidad de Financiación de los Organismos de 0.22 por ciento del PIB. Por lo demás. llama la atención el déficit registrado por la Administración Local. Ha sido debido. principalmente, a la liquidación definitiva de los recursos del sistema de financiación de 2020. De cara al futuro se espera que las administraciones progresen en la consolidación fiscal. Este proceso será más lento en la Administración Central. básicamente por las transferencias que deberá seguir haciendo a la Seguridad Social, como veremos más adelante.


\subsection{Ingresos y gastos de los su bs·ectores de las AAPP}

A continuación. nos referiremos al reparto de los ingresos y los gas1os entre los di!:>tintos niveles de las AAPP en 2022. En este caso, en los datos de la tabla 9 prescindiremos de la clasificación entre operaciones co1Tientes y de capital, por no ser relevante.

\textbf{TABLA DE INGRESOS Y GASTOS}

1 . El mayor tamaño, tanto en gastos (50,2\%) como en ingresos. de la Administración Central

2. ! lay una nueva rúbrica ··Transferencias entre AAPP" que se cancela en el agregado. En gran parte se debe a las transferencias debidas al sistema de financiación autonómico y local. y a la provisión de financiación de la Administración Central a las administraciones territoriales para necesidades que surgen anualmente. 1.'.0mo por ejemplo en Sanidad. en Educación y en Dependencia. También se incluyen las transferencias del Estado a la Seguridad Social para la financiación de sus gastos, incluyendo los ·'gastos impropios'': y al SEPE para gastos de desempleo y políticas no contributivas de empleo. La magnitud de esta rúbrica, 1 2.5 por ciento del PIB (la mayor del gasto de la AC), pone de relieve el papel central del Estado dentro de las AAPP. Dada la naturaleza de las transferencias aquí incluidas, la magnitud de esta rúbrica puede variar sensiblemente de año en año.

\begin{specialblock}{\textbf{Gastos impropios:} ¿Qué son?}
    Conviene detenernos en el tema de los '·gastos impropios ... por c;;u conexión con el crecimiento de los gastos con-ientes y el riesgo para la sostenibilidad de las finanzas de las AAPP que mencionamos al p1incipio del lema.

    El principio de separación de fuentes. que el Pacto de Toledo ha incluido entre sus recomendaciones desde su inicio en 1995, exige que las pensiones y otras prestaciones contributivas de la Seguridad Social ,;e financien básicamente con cotizaciones sociales, y que lo�� complememos de mínimo:, ) prestaciones asistenciales se financien con impuestos generales que recauda el Estado ) éste transfü:rc a la Seguridad Socia1.\footnote{%
        De esta forma se busca separar la financiación de las prestaciones que responden a un principio de solidaridad. asegurando unas rentas mínimas a aquellas personas que carecen de otros recursos. de las que funcionan más como un seguro. ofreciendo pagos proporcionales a las cuotas satisfechas si ��e producen determinadas contingencias (enfermedad. jubilación, viudedad, maternidad o paternidad ... ).
    } El componente contributivo o de seguro de la Seguridad Social debería, en principio, autotinanciarse en circunstancias normales. pero no sería jw,to cargarle además el coste de prestaciones externas que siguen una lógica distinta.

    El problema es que. cada vez más. las cotizaciones sociales han tenido que complementarse con aportaciones directas del Esiado, por ejemplo, para complementar las pensiones contributivas más bajas hasta una cuantía mínima (complemento de mínimos). Así mismo, en la recomendación I del Informe de la Comisión del Pacto de Toledo40 . se afirma que el déficit de la Seguridad Social se debe fundamentalmente a que el sistema sigue haciéndose cargo Je una serie de ··gastos impropios .. que. _por no tratarse de prestaciones contributivas. deberían financiarse con transferencias del Estado en vez de hacerlo con cotizaciones sociales. Diversos expertos discuten algunos de los gastos impropios que cita el lnfonne, por considerarlos dentro del ámbito contributivo. Tales son los casos de la prestación contributiva por nacimiento y cuidado de hijos y los costes de funcionamiento41.

    Controversias aparte, los gastos de la Seguridad Social plantean un problema creciente para las finanzas de las AAPP, que se revela con claridad al analizar las relaciones entre la AC y los FFSS. La reforma del sistema público de pensiones materializada en la Ley 21/2021 y en el ROL 2/202 3. que introdujo refotmas tanto en las pensiones de jubilación como en las cotizaciones. tiene la intención de. al menos. atenuar este problema. Para un análisis más profundo del sistema español de pensiones, puede consultarse De la Fuente (2023b). aunque el análisis que se lleva a cabo reviste alguna complejidad.
\end{specialblock}


3. La A R tiene mayor peso en Remuneración de asalariados (58 por ciento del total) y en Consumos inte1medios. dadas las competencias asumidas en Sanidad. Educación y Dependencia.

4. En materia de Formación Bruta de Capital, el gasto está bastante equitativamente repartido entre la AC y cada una de las dos las administraciones territo1iales.

5. El reparto de impuestos y cotizaciones sociales se corresponde con lo visto al analizar el reparto de la presión fiscal, aunque el desglose no es el mismo. Aquí. sin embargo, la lGAE atribuye a la AR su parte de la recaudación de los impuestos especiales. a pesar de ser un impuesto estatal y de controlar el Estado los tipos impositivos. Aquí. debe tenerse en cuenta que la Comunidad Autónoma del País Vasco y su:, Diputaciones Forales. además de la Comunidad Autónoma de Nava1Ta, tienen atribuidos la totalidad de los impuestos recaudados en sus territorios.


\begin{specialblock}{ANÁLISIS DE LAS CUENTAS DE LAS CCAA}
    Este tema tiene un contenido fundamentalmente descriptivo. con breves explicaciones de los factores subyacentes a los datos.
    
    No obstante. hay centros de estudios, como f-EDEA. que suelen publicar trabajos de análisis como el citado en la bibliografía: Oc la rucnte (2023). Este estudio presenta un análisis del desempelio de las CCAA en 2022. Entre sus conclusiones destacan las siguientes:
    
    1. Desde 2003 los gastos y los ingresos han sido fuertemente pro-cíclicos 
    
    2. Preocupa el alto nivel de deuda de las CCAA
    
    3. Preocupa, asimismo, que la mejora del saldo presupuestario se apoya en parte en factores anómalos como una inversión atípicamente baja y fuertes subvenciones a los intereses a través de los mecanismos estatales de financiación y liquidez

    4. El repunte del déficit en 2022 debe bastante a las circunstancias atípicas del ��jercicio y no se extiende al déficit subyacente. que de hecho se reduce casi la mitad. Sin embargo. preocupa que, de momento, los gastos e ingresos subyacentes no parecen volver a los niveles preandemia.
\end{specialblock}

\subsection{Las cuentas financieras de los su bscctores de las AAPP}

Lo más destacable de las cuentas financieras de los subsectores de las AAPP es que ponen todavía más de relieve el papel de la Administración Central dentro del conjunto de las AAPP y de la economía española. A diferencia de lo que hicimos al principio del tema al repasar las cuentas consolidadas de las AAPP. en vez de referimos a las variaciones de las magnitudes, nos referiremos a los niveles por tener más interés.

En primer lugar. hay que destacar que la AC incluye 110 sólo al Estado y entidades de"pendientes no productoras de mercado.· sino también fondos creados para determinados fines. Tales son los casos del Fondo para la Adquisición de Activos Financieros (FAAF). del Fondo de Reestructuración Ordenada Bancaria (FROB). del Fondo de Amortización del Déficit Eléctrico (fADE). del Fondo de Financiación a las Comunidades Autónomas FFCCAA (creado en 2014. que asumió las deudas del Fondo de Liquidez Autonómico (FLA) y del Fondo Financiero de los Pagos a Pro\'eedores (FFPP)), y del fondo de Garantía de Depósitos (FGD). A esto!> fondos hay que aiiadir la Sociedad de Gestión de Activos Procedentes de la Reestructuración Bancaria (SAREB). Así mismo, dentro del Estado. se encuadra el préstamo del MEDE.

En segundo lugar, una rar1e de los activos financieros de la Administración Central los tiene frente a la AR, la AL y lo!> FFSS. En el caso de las CCAA. a travt'S del FFCCA, y en el caso de los FFSS. por medio del pr0stamo a la Seguridad Social.

En concreLO, para una deuda total tsin consolidar) PDE de la AC, a finales de 2022. del l 02.4 por ciento del Pll3. 22,5 puntos de PIB se deben a la necesidad de endeudarse para proporcionar financiación a la AR. la AL y los ffSS4��. En concreto. el préstamo del Cstado a la Seguri<lad Social asciende a 8 puntos del PIB4\ mientras que la deuda de la AR y la AL frente al Estado es de 14.5 puntos del PIB (14 puntos de la AR y 0 ,5 puntos de la AL44).

Y, por lo que se re ti ere a los activos financieros de la J\C relacionados con el sector financiero y eléctrico. cabe decir que, de los 102.4 puntos de PIB de deuda, 100,3 corresponden al Estado: de los cuales. 1.5 punto�� corresponden al préstamo del MEDE. Dentro de las otras unidades de la AC, 2,4 puntos corresponden a la SAREB. 0,6 puntos corresponden al F ADE y O, 1 puntos al FR0B45.

Por último. la distribución del stock de deuda PDE consolidado4�� de las AAPP en 2022, del 113,2 por ciento del PIB entre Administraciones fue:

AC: 79.8
AR: 23,9
AL: 1.7
FFSS: 7,8

En 2022. las autonomías menos endeudadas. en términos de sus P!Bs, eran: Madrid, Canarias y las comunidades forales (entre el 13 y el 14 por ciento del Pll3). mientras que las más endeudadas eran: Murcia. Castilla-la Mancha. Cataluña (entre el 32 y el 34 por ciento del PlB) y, sobre todo, la Comunidad Valenciana (44A por ciento del PlB).

Por último, cabría referirse a los pasivos contingentes de las AAPP y a la deuda de las empresas públicas no encuadradas dentro de las AAPP. No ohstante. al no ser obligatorias estas referencias. las relegamos a los anexos 2 y 3.
 




\clearpage
\section{Finanzas de las Administraciones Públicas: Marco jurídico-conceptual}
\puntosclave{
    
}

Cerramos esta serie de tres temas con el que, en cierto sentido, les da su sentido de conjunto. El Tema 21 nos enseñó el marco jurídico y el ciclo del presupuesto del Estado; el Tema 22 nos enseñó a leer ese mismo presupuesto —y el de sus organismos dependientes— con las tres grandes clasificaciones de la contabilidad presupuestaria. Este último tema cambia, deliberadamente, tanto el objeto como el lenguaje: ya no hablamos solo del Estado y sus entes limitativos, sino del conjunto de las Administraciones Públicas españolas —Estado, Comunidades Autónomas, Entidades Locales y Seguridad Social—; y ya no hablamos en el lenguaje de la contabilidad presupuestaria española, sino en el lenguaje internacional de la contabilidad nacional. Este primer bloque construye, antes que ninguna cifra, el porqué de ese doble cambio: por qué tiene sentido tratar a un conjunto tan heterogéneo de Administraciones como si fueran una sola cosa, por qué hace falta aislar las operaciones financieras del resto, qué son exactamente los "subsectores" que vamos a usar durante todo el tema, y solo entonces, con ese armazón ya construido, la fotografía actual de las cuentas públicas españolas y su trayectoria prevista.

\subsection{Las AAPP como única unidad económica}
El problema de partida: ¿qué sentido tiene tratar como "una sola cosa" a miles de entidades jurídicamente distintas?

[Probable] Empiezo por la pregunta que cualquier lector atento debería hacerse antes de aceptar, sin más, la premisa de todo este tema. España tiene un Estado, diecisiete Comunidades Autónomas y dos Ciudades Autónomas, más de ocho mil municipios, decenas de Diputaciones Provinciales, y un sistema de Seguridad Social con sus propias entidades gestoras —miles de personas jurídicas distintas, cada una con su propio presupuesto, su propio Parlamento o Pleno al que responde, su propia contabilidad—. ¿Qué sentido tiene, entonces, sumar todo esto y hablar de "las finanzas de las AAPP" como si fueran, efectivamente, una única unidad económica con un único balance y una única cuenta de resultados? [Seguro] La respuesta no es una convención arbitraria de los estadísticos: es la aplicación, a su escala más amplia posible, de exactamente la misma técnica de consolidación que ya fundamentamos con todo detalle en el Tema 22 —recuerda la distinción entre consolidación intrasectorial (dentro de un mismo subsector, el ejercicio que hicimos en aquel tema) e intersectorial (entre distintos subsectores, el ejercicio que hacemos aquí)—. Este tema es, literalmente, el segundo y último escalón de un proceso de consolidación que empezamos a construir hace ya dos temas: primero consolidamos el Estado con sus organismos dependientes; ahora consolidamos ese resultado con las Comunidades Autónomas, las Entidades Locales y la Seguridad Social, para obtener, por fin, la magnitud agregada que tiene sentido comparar con el PIB, con otros países, y con los umbrales de la propia gobernanza fiscal europea que ya trabajamos en el Tema 21.

La sectorización: una decisión de diseño estadístico con su propia genealogía

[Seguro] Este ejercicio de agrupamiento no es exclusivo de las Administraciones Públicas: es una pieza central de toda la arquitectura del Sistema de Cuentas Nacionales que ya conocemos, desde el Tema 22, por su vínculo directo con Richard Stone. El SCN divide la totalidad de la actividad económica de un país en un número reducido de sectores institucionales —hogares, sociedades no financieras, sociedades financieras, Administraciones Públicas, e Instituciones sin Fines de Lucro al Servicio de los Hogares—, más el resto del mundo como sector exterior, precisamente para poder analizar la economía como un sistema de flujos circulares entre un número manejable de actores, en lugar de como una masa indiferenciada de millones de agentes individuales. [Probable] Las Administraciones Públicas —el Sector S.13 del SEC-2010— son, dentro de esta arquitectura, el sector que agrupa a todas las unidades institucionales que cumplen el doble criterio que ya trabajamos, sin desarrollarlo del todo, en el Tema 22: estar sometidas a control público efectivo, y ser productoras no de mercado —proporcionar bienes y servicios de forma gratuita o a precios que no cubren sus costes de producción de forma "económicamente significativa"—. [Seguro] Este criterio doble es, precisamente, lo que permite tratar como "una sola unidad" a entidades tan dispares como el Ministerio de Defensa y un pequeño ayuntamiento rural: ambos comparten la misma naturaleza económica de fondo —no venden lo que producen en un mercado real—, aunque no compartan absolutamente nada en términos de tamaño, competencias o nivel de gobierno.

El problema que esta "unidad económica" plantea, y que ningún tema anterior de esta serie ha tenido que resolver: ¿qué produce exactamente un sector que no vende nada?

[Seguro] Aquí llegamos al problema teórico más profundo y menos conocido de todo este epígrafe, y que conviene desarrollar con el detalle que merece porque revela una fragilidad metodológica real en el corazón mismo de la contabilidad nacional. Si una empresa privada vende sus productos en el mercado, medir su producción es sencillo: se observa el precio al que vende, se multiplica por la cantidad vendida, y se obtiene su valor de producción. Pero, ¿cómo se mide la "producción" del sector público, cuando la inmensa mayoría de lo que produce —defensa nacional, justicia, seguridad ciudadana, administración general— no tiene precio de mercado observable en absoluto, porque nunca se vende a nadie? [Seguro] La convención que la contabilidad nacional ha empleado tradicionalmente para resolver este problema —vigente en el Reino Unido, según reconoce la propia literatura especializada, desde los años sesenta hasta 1998, y todavía dominante en gran parte del mundo— es la llamada convención "output = input": el valor de la producción no de mercado del sector público se define, por pura convención contable, como exactamente igual al coste de los factores empleados para producirla —remuneración de asalariados, consumos intermedios, y el consumo de capital fijo—. [Seguro] Esta convención tiene una implicación que, formulada con precisión, resulta genuinamente sorprendente y merece retenerse como el gancho conceptual de todo este epígrafe: si el valor de lo que el sector público "produce" se define, por definición, como idéntico a lo que cuesta producirlo, entonces la productividad del sector público —la relación entre lo que se produce y lo que se gasta en producirlo— es, por pura construcción matemática, siempre constante, con una tasa de crecimiento exactamente igual a cero. No es que la productividad del sector público sea difícil de mejorar: es que, bajo esta convención contable, es literalmente imposible que aparezca mejorando, con independencia de que un hospital atienda a más pacientes con los mismos recursos, o de que una escuela eduque mejor a los mismos alumnos, porque el propio método de medición excluye, por diseño, esa posibilidad.

[Seguro] Esta fragilidad metodológica no pasó desapercibida, y su corrección tiene nombre y fecha propios: el National Statistician británico, Len Cook, encargó en diciembre de 2003 al economista Sir Anthony Atkinson —Warden del Nuffield College de Oxford, y la misma figura, no por casualidad, cuyo teorema de Atkinson-Stiglitz ya trabajamos con todo detalle en el Tema 13 a propósito de la diferenciación de tipos de IVA— una revisión independiente de la metodología de medición de la producción y la productividad del sector público en la contabilidad nacional. El Informe Atkinson, publicado en 2005, estableció un conjunto de principios que siguen citándose hoy como referencia internacional —entre ellos, que la medición del output no de mercado debe seguir, en la medida de lo posible, un procedimiento paralelo al empleado para la producción de mercado, y que esa producción debe medirse directamente —contando operaciones quirúrgicas realizadas, alumnos efectivamente educados, ajustados por calidad— en lugar de inferirse mecánicamente del coste de los factores empleados—. [Probable] La magnitud de lo que está en juego en este debate, aparentemente técnico, es enorme: la producción no de mercado del sector público representa, según las propias cifras británicas recientes, en torno al 20% del PIB del Reino Unido, con cifras similares —entre el 19% y el 24%— en el resto de los países del G7 salvo Estados Unidos —donde ronda solo el 14%, reflejo de un Estado del bienestar de menor tamaño relativo—; ignorar o medir mal la productividad de una quinta parte de toda la economía de un país no es, ni mucho menos, un detalle técnico menor, sino una laguna que puede distorsionar de forma sustancial cualquier análisis agregado de la productividad de un país entero. [Probable] El propio Informe Atkinson influyó directamente en la revisión del Sistema de Cuentas Nacionales de 2008 (SCN 2008), que recogió sus principios como recomendación general para todos los países que siguen este estándar internacional —aunque, conviene señalarlo con la misma honestidad con la que hemos tratado cada debate de este proyecto, la aplicación práctica de la medición directa de la productividad pública sigue siendo, más de dos décadas después del informe, un terreno de desarrollo metodológico activo y desigual entre países, con revisiones posteriores —como la del propio Reino Unido, con el llamado Bean Review de 2016 y una revisión independiente más reciente encargada en 2023— que siguen intentando perfeccionar un problema que el propio Atkinson advirtió, ya en 2005, que no tenía una solución sencilla ni definitiva.

Por qué, pese a esta fragilidad, la ficción de la "unidad económica" sigue siendo indispensable

[Probable] Conviene cerrar este epígrafe explicando por qué, pese al problema de medición que acabamos de desarrollar con tanto detalle, ningún país ni ninguna institución internacional ha abandonado nunca la idea de tratar a las Administraciones Públicas como una única unidad económica coherente. La razón de fondo es la misma que ya intuimos en la introducción: sin esta ficción metodológica, sería imposible aplicar cualquiera de las reglas fiscales que trabajamos con tanto detalle en el Tema 21 —el déficit del 3%, la deuda del 60%, el propio art. 135 CE—, porque todas ellas exigen, como condición técnica previa e insoslayable, una única cifra agregada contra la que medir el cumplimiento, no una colección dispersa de miles de balances municipales, autonómicos y estatales sin ningún criterio común de agregación. La ficción de la unidad económica no es, por tanto, un capricho académico de los estadísticos: es la precondición técnica de toda la arquitectura de disciplina fiscal, nacional y europea, que hemos ido construyendo a lo largo de los tres temas de esta serie — el mismo tipo de infraestructura conceptual silenciosa, y rara vez explicada con el detalle que merece, que ya identificamos, en otro contexto, al hablar de por qué la clasificación presupuestaria es el lenguaje necesario del principio de especificación en el Tema 22.


\subsection{Por qué se aíslan las operaciones financieras: el fundamento teórico}
Como sabemos, las operaciones financieras no computan en el saldo público, mientras que las operaciones no financieras sí lo hacen. Pero, ¿por qué existe esta frontera y qué principio económico de fondo la sostiene? Es la pregunta que este epígrafe responde, porque es, en última instancia, la pregunta que separa un déficit público real de un déficit público maquillado.

### La secuencia de cuentas del SCN: el deficit nace en la cuenta de capital, no en la financiera

[Seguro] El Sistema de Cuentas Nacionales organiza toda la actividad económica de un sector en una secuencia ordenada de cuentas, cada una alimentando a la siguiente. La cuenta de producción da lugar a un excedente; la cuenta de distribución y utilización de la renta —donde entran salarios, impuestos, cotizaciones, transferencias— desemboca en el ahorro del sector, la misma magnitud que ya trabajamos en el Tema 23 original; la cuenta de capital —donde se registra la formación bruta de capital y las transferencias de capital— desemboca, por fin, en la magnitud que estábamos buscando: la capacidad o necesidad de financiación, es decir, el propio déficit o superávit público. [Seguro] Es solo después de esta secuencia, y como su consecuencia lógica y no como una fuente independiente de resultado, cuando entra en juego la cuenta financiera: registra, exclusivamente, cómo se ha financiado esa necesidad de financiación ya determinada —mediante qué instrumentos concretos: emisión de deuda, disposición de depósitos, concesión o cancelación de préstamos—. [Probable] La secuencia importa, y mucho, porque revela algo que conviene fijar con toda precisión: la cuenta financiera no genera déficit ni superávit; se limita a reflejar, como un espejo, la decisión de financiación de un resultado que ya estaba determinado, de forma completa y cerrada, en la cuenta de capital que la precede. El déficit público es, en este sentido estricto, un concepto que pertenece por completo al ámbito de las operaciones reales —consumo, inversión, transferencias sin contrapartida—, nunca al de las operaciones financieras que simplemente instrumentan su cobertura.

### Patrimonio neto: la magnitud que la distinción protege

[Seguro] El concepto que da sentido económico a toda esta arquitectura es el de patrimonio neto —activos menos pasivos—. Una operación no financiera —pagar un salario, conceder una transferencia sin contrapartida, comprar un bien de consumo— reduce el patrimonio neto de quien la realiza, de forma efectiva y sin ningún retorno compensatorio: el dinero sale, y no se recibe a cambio ningún activo de valor equivalente. Una operación financiera —comprar deuda pública, conceder un préstamo, ingresar el reembolso de un crédito anterior— es, en cambio, un puro intercambio: el efectivo que sale se sustituye, en el balance, por un activo financiero de idéntico valor —un derecho de cobro—, de modo que el patrimonio neto del sector permanece, en el instante mismo de la operación, exactamente igual que antes de realizarla. [Probable] Esta es la razón económica de fondo, y no una convención arbitraria de los estadísticos, por la que el déficit público —la magnitud que de verdad mide si el sector público se está empobreciendo o enriqueciendo netamente en un ejercicio— excluye, por definición, cualquier operación que sea puro intercambio de activos de igual valor: si las incluyéramos, estaríamos confundiendo un cambio en la composición de la riqueza pública —efectivo por un préstamo, por ejemplo— con un cambio en su magnitud total, que es, precisamente, lo único que el déficit pretende capturar.

### Lo que ocurre cuando esta frontera se explota deliberadamente: el caso griego

[Seguro] La importancia práctica de esta distinción, lejos de ser un tecnicismo contable sin consecuencias, quedó demostrada de la forma más dramática posible con el episodio del fraude estadístico griego, descubierto entre 2009 y 2010 y que constituye, probablemente, el caso mejor documentado de toda la historia de la gobernanza fiscal europea sobre lo que ocurre cuando un Estado explota deliberadamente la frontera entre operación financiera y no financiera para maquillar sus cifras. [Seguro] A comienzos de 2002, poco después de la adopción del euro por Grecia, el Tesoro griego —bajo la dirección de Christoforos Sardelis, responsable de la Agencia griega de Gestión de la Deuda Pública— concertó con Goldman Sachs una operación de permuta de divisas cruzadas (cross-currency swap): deuda griega denominada en dólares y en yenes se intercambiaba por deuda denominada en euros, utilizando, para fijar los tipos de cambio de esa permuta, cifras deliberadamente distintas de las cifras reales de mercado —los llamados "tipos de cambio ficticios"—, de modo que la operación permitía a Grecia obtener, de facto, una inyección adicional de liquidez de aproximadamente 2.800 millones de euros, sin que esa cantidad apareciera contabilizada como nueva deuda en las estadísticas remitidas a Eurostat, precisamente porque una permuta de divisas se clasifica, en las normas estadísticas de la época, como una operación financiera —un mero intercambio de instrumentos de igual valor nominal—, no como el endeudamiento adicional que, en términos económicos reales, efectivamente era. [Seguro] Cuando la permuta llegó a su vencimiento, años después, la propia deuda oculta —que para entonces había crecido hasta unos 5.700 millones de euros por el coste acumulado de la operación— tuvo que aflorar en las cuentas públicas griegas, contribuyendo a la crisis de credibilidad estadística que precedería, apenas unos años después, al colapso de la confianza de los mercados en la deuda soberana griega y al inicio de la crisis del euro. [Seguro] Eurostat, en su revisión de septiembre de 2004, ya había tenido que corregir al alza el déficit griego declarado de 2002, del 1,2% inicialmente comunicado hasta el 3,7% real —una revisión que, con datos posteriores, llegaría hasta el 5,2%—, y el propio Goldman Sachs reconocería públicamente que este tipo de operaciones eran "una de varias técnicas que muchos gobiernos europeos utilizaron para cumplir los términos del Tratado" de Maastricht —confirmando que el caso griego, aunque el más extremo y mejor documentado, no fue un episodio aislado: Italia recurrió a técnicas de titulización similares, con la ayuda de otro banco estadounidense, para maquillar sus propias cifras de deuda en el mismo periodo.

### La respuesta regulatoria: por qué Eurostat ha ido cerrando, uno a uno, los resquicios de la "contabilidad creativa"

[Probable] Este episodio, y otros de naturaleza similar detectados en distintos Estados miembros a lo largo de los años —titulizaciones de ingresos futuros que trasladaban deuda fuera del balance público, operaciones de venta y posterior arrendamiento (sale and leaseback) de activos públicos, y muy señaladamente la clasificación de las colaboraciones público-privadas (PPP), donde el propio riesgo de construcción y de disponibilidad de una infraestructura puede, según cómo se estructure contractualmente, mantenerse artificialmente fuera del balance de las Administraciones Públicas aunque el Estado siga siendo, en la práctica económica real, quien soporta la mayor parte del riesgo—, explican por qué Eurostat ha ido desarrollando, década tras década, un conjunto de normas cada vez más detalladas y más restrictivas sobre qué operaciones pueden clasificarse como puramente financieras y cuáles deben reclasificarse, a efectos estadísticos, como operaciones reales con impacto directo en déficit —el llamado Manual on Government Deficit and Debt, actualizado de forma continua precisamente para cerrar, uno a uno, los resquicios que la práctica de la "contabilidad creativa" ha ido descubriendo—. [Suposición] Es, en definitiva, el mismo patrón de fondo que ya hemos visto repetirse en otros contextos de este proyecto —recuerda el propio Pacto de Estabilidad y Crecimiento, cuya "falta de credibilidad sancionadora" ya señalamos en el Tema 21—: una regla fiscal numérica y rígida genera, casi por necesidad, un incentivo a explotar cualquier margen técnico de definición que separe lo que cuenta de lo que no cuenta, y la distinción entre operación financiera y no financiera que hemos fundamentado en este epígrafe es, precisamente, uno de los terrenos donde esa tensión entre la letra de la norma y su finalidad económica real se ha manifestado con mayor claridad —y con mayores consecuencias— en toda la historia reciente de la gobernanza fiscal europea.

\subsection{Los subsectores: qué son, por qué existen, qué distingue a cada uno}
El SEC-2010 subdivide el Sector S.13 —Administraciones Públicas— en cuatro subsectores con codificación propia: S.1311 (Administración Central), S.1312 (Comunidades Autónomas, denominadas en la terminología general del sistema "administraciones de los estados federados" para adaptarse a países de estructura federal), S.1313 (Corporaciones Locales) y S.1314 (Fondos de la Seguridad Social). [Probable] Pero limitarse a citar esta codificación sería quedarse en la superficie del problema: la pregunta genuinamente interesante no es cómo se llama cada subsector, sino por qué el sistema estadístico internacional considera necesario, y no meramente decorativo, dividir lo que en el epígrafe 1.1 acabamos de fundamentar como "una única unidad económica" en cuatro piezas distintas y analizadas por separado.

### Por qué subdividir: tres razones que van más allá de la mera comodidad expositiva

[Probable] La primera razón ya la conocemos, porque la trabajamos con todo detalle en el Tema 21: las reglas fiscales de la LOEP —déficit, deuda, regla de gasto— se aplican de forma diferenciada por subsector, no de manera uniforme al conjunto; sin la subdivisión estadística, sería sencillamente imposible verificar si cada uno de esos límites diferenciados se cumple. La segunda razón es de naturaleza institucional: cada subsector responde, como ya vimos, ante un órgano de control democrático distinto —Cortes Generales, Parlamentos autonómicos, Plenos municipales—, con grados de autonomía normativa y de responsabilidad política radicalmente distintos entre sí, de modo que agregarlos sin distinción ocultaría diferencias de comportamiento fiscal que resultan, en la práctica, decisivas. [Seguro] La tercera razón, la que menos se suele explicitar y la que más interesa para el propósito de este epígrafe, es de naturaleza comparativa internacional: la subdivisión en subsectores es, precisamente, lo que permite medir el grado de descentralización fiscal de un país frente a otro —qué proporción del gasto público total gestiona el nivel central frente a los niveles territoriales—, un indicador que organismos como la OCDE y el FMI publican de forma sistemática y que resulta, aplicado al caso español, genuinamente revelador.

### España, el caso que no encaja en ninguna casilla: ni país federal, ni país unitario

[Seguro] Y aquí llega el dato que merece el desarrollo más extenso de todo este epígrafe, porque revela algo que la simple enumeración de los cuatro subsectores nunca dejaría ver por sí sola. La propia OCDE, en sus estadísticas comparadas de descentralización fiscal, no incluye a España ni en la categoría de "países federales" —Alemania, Estados Unidos, Canadá, Austria, Suiza— ni en la de "países unitarios" convencionales: la sitúa en una tercera categoría, la de "países regionales", que España comparte, según reconoce la propia OCDE, únicamente con Colombia, por tratarse de países con "una estructura política muy descentralizada, con una autonomía importante de sus entidades territoriales", sin llegar a ser, en sentido constitucional estricto, Estados federales. [Seguro] Las cifras que sostienen esta clasificación son, además, notablemente elocuentes: la recaudación tributaria gestionada por las administraciones autonómicas españolas pasó del 5% del total en 1995 al 15,7% en 2021 —una trayectoria de descentralización fiscal que la propia OCDE describe como la más intensa de todos los países del área desde mediados de los años noventa—, situando a España, en términos de recaudación regional, muy por delante de países federales europeos como Austria (2%) y a un nivel comparable al de Alemania (24,7%) o Suiza (24,5%), aunque todavía por detrás de Canadá (39,6%), el país con mayor descentralización tributaria regional de toda la OCDE. [Probable] En términos de gasto —no ya de ingreso— la asimetría es todavía más marcada: España es, según los estudios comparados disponibles, el único país europeo con gobiernos intermedios donde el gasto regional en proporción del PIB iguala al de la propia Administración Central, con más del 90% del gasto público en sanidad y en educación gestionado por las Comunidades Autónomas —un grado de descentralización de estas dos grandes políticas del Estado del bienestar que no tiene parangón entre los grandes países federales de nuestro entorno—.

[Suposición] Este dato, que a primera vista parece una simple curiosidad estadística, tiene una implicación de fondo que conviene explicitar como cierre de esta comparación: España ha construido, a lo largo de más de cuatro décadas desde la Constitución de 1978 y la LOFCA de 1980, un grado de descentralización fiscal funcionalmente equivalente al de un Estado federal —en volumen de recursos gestionados, en autonomía tributaria efectiva, en responsabilidad política sobre las grandes políticas del bienestar—, sin haber adoptado nunca, formalmente, una arquitectura constitucional federal. Es el mismo fenómeno que la ciencia política española denomina, con frecuencia, "federalización sin federalismo" o, en la terminología más clásica del propio texto constitucional, el "Estado de las Autonomías" del art. 2 CE —un modelo sui generis que ni la propia Constitución de 1978 definió con precisión terminológica, dejando a la práctica política y a la jurisprudencia constitucional posterior la tarea de ir perfilando, caso por caso y durante décadas, el alcance real de esa descentralización—.

### Qué distingue a cada subsector, más allá del catálogo de entidades

[Seguro] Con esta comparación ya fijada, recorro brevemente la naturaleza distintiva de cada uno de los cuatro subsectores, no como mero catálogo de entidades —eso ya lo hicimos, con mucho detalle, en el Tema 22 para la Administración Central— sino como categorías analíticas con un comportamiento fiscal propio y diferenciado. La Administración Central es la única con capacidad efectiva de actuar como prestamista de última instancia frente a los demás subsectores —recuerda el préstamo del Estado a la Seguridad Social, o el Fondo de Liquidez Autonómica que ya trabajamos—, y la que concentra, coherentemente con ese papel, el mayor límite de deuda dentro del reparto LOEP (44% del PIB). La Administración Regional es la que mejor ilustra la asimetría constitucional española que acabamos de cuantificar —con el propio contraste interno entre el régimen común y el régimen foral que ya trabajamos en el Tema 21—. La Administración Local, sujeta al criterio más restrictivo de mera "suficiencia financiera" y no de "autonomía financiera" plena, es la que, de forma sistemática desde 2012, muestra superávits estructurales recurrentes —el reflejo exacto, en el terreno del comportamiento fiscal observado, de la menor autonomía normativa que la propia Constitución le reconoció—.

### Los Fondos de la Seguridad Social: el subsector que no es, en realidad, un "nivel" de gobierno

[Probable] Cierro con la pieza que merece la reflexión conceptual más cuidadosa de las cuatro, porque su inclusión en esta lista genera, a menudo, una confusión que conviene resolver explícitamente: a diferencia de los otros tres subsectores, organizados según un criterio territorial —central, regional, local—, los Fondos de la Seguridad Social se organizan según un criterio completamente distinto: funcional. No son un "nivel" de gobierno que se sitúe, en la jerarquía territorial, entre la Administración Regional y la Local, ni por encima ni por debajo de ninguna de ellas: son, más bien, un corte transversal de toda la actividad pública española dedicada específicamente a la función de aseguramiento social —pensiones, desempleo, dependencia—, que opera, con sus propias entidades gestoras, en todo el territorio nacional simultáneamente, sin distinción territorial alguna, salvo la peculiaridad ya conocida del régimen foral. [Suposición] Esta diferencia de naturaleza —tres subsectores organizados por territorio, uno organizado por función— no es un simple detalle taxonómico: explica, de hecho, por qué la Seguridad Social puede convivir, sin contradicción lógica alguna, con cualquiera de los otros tres niveles territoriales sin competir nunca con ellos por el mismo espacio competencial, y por qué, precisamente por esta naturaleza funcional y no territorial, resulta el subsector más sencillo de tratar como una unidad estadística coherente y homogénea en toda España, frente a la enorme heterogeneidad interna que sí caracteriza, como ya hemos visto, tanto al nivel regional como al local.


## 1.4. Situación actual y perspectivas: agregada y por subsector

### Cuadro-resumen (verificar y actualizar contra el Informe trimestral de la IGAE vigente antes de usar en examen)

| **% del PIB** | 2019 | 2020 | 2021 | 2022 | 2023(p) | 2024(p) |
|---|---|---|---|---|---|---|
| Ingresos totales | 39,2 | 41,8 | 43,7 | 43,0 | 43,4 | 43,3 |
| Gastos totales | 42,3 | 51,9 | 50,6 | 47,8 | 47,3 | 46,3 |
| Cap./Nec. de financiación | -3,1 | -10,1 | -6,9 | -4,8 | -3,9 | -3,0 |
| PIB nominal (mm €) | 1.245,5 | 1.118,0 | 1.206,8 | 1.327,1 | 1.408,1 | 1.494,0 |

### Los grandes resultados de 2022, y la advertencia de sostenibilidad que ningún cuadro de cifras estáticas puede transmitir por sí solo

**[ICEX-CECO]** El gasto total alcanzó el 47,8% del PIB en 2022 —5,5 puntos por encima de 2019—, sin previsión de retorno a los niveles previos a la pandemia; los ingresos, el 43,0% —3,8 puntos por encima—; la necesidad de financiación, el 4,8%, con una senda descendente prevista hasta el 2,5% en 2025-2026. [Seguro] Pero la cifra que merece el mayor desarrollo, y que el documento de referencia menciona sin desarrollarla con el aparato que hoy tenemos a nuestra disposición, es la del **gasto corriente per cápita**: +26,4% en términos reales entre 2008 y 2023, frente a un crecimiento del PIB real per cápita de apenas el **3,0%** en el mismo periodo —una divergencia de casi nueve veces entre ambas magnitudes que, sostenida en el tiempo, plantea una pregunta de sostenibilidad que ningún ejercicio presupuestario anual, por sí solo, puede resolver—. [Seguro] La confirmación más reciente y más autorizada de esta preocupación llega del propio **Informe de Envejecimiento de 2024** de la Comisión Europea (*Ageing Report*), elaborado cada tres años por el Grupo de Trabajo sobre Envejecimiento y Sostenibilidad creado en 1999 para asesorar a la Comisión: España aparece, en esta edición, como el **país de la Unión Europea donde más aumentará el gasto en pensiones** de toda la proyección hasta 2070 —un incremento de **3,3 puntos de PIB en 2050 y de 5 puntos en 2070**, con una revisión al alza de **6,5 puntos** respecto a la propia edición de 2021, la mayor corrección de todos los Estados miembros—. [Probable] Lo verdaderamente revelador de este dato es su propia causa: la reforma de pensiones que ya trabajamos con detalle en el Tema 22 —revalorización por IPC, supresión del factor de sostenibilidad, Mecanismo de Equidad Intergeneracional— es, precisamente, la responsable directa de este empeoramiento de la proyección a largo plazo respecto al escenario anterior: al restaurar la garantía de poder adquisitivo y eliminar el mecanismo que automáticamente habría contenido el gasto futuro, la propia reforma diseñada para dar más seguridad a los pensionistas actuales ha encarecido, de forma medible y cuantificada por la propia Comisión, la factura que las próximas décadas tendrán que financiar. [Seguro] Y aquí aparece un dato de gobernanza que conecta directamente con lo que ya construimos en el Tema 22: la propia regulación del MEI establece que, si a partir de **2033** el gasto real en pensiones se desvía de la proyección de este mismo Informe de Envejecimiento de 2024 —usado explícitamente como referencia normativa—, se activará el nuevo Fondo de Reserva con un límite de disposición anual del 0,2% del PIB, y si esto resultara insuficiente, el Gobierno deberá negociar con los agentes sociales, en el marco del Pacto de Toledo, nuevas medidas correctoras —de reducción del gasto o de incremento de cotizaciones—. Es la primera vez, en toda esta serie de temas, que vemos un informe técnico de proyección convertirse, por diseño normativo explícito, en el propio **disparador automático** de una futura reforma fiscal.

### Las operaciones financieras y el ajuste déficit-deuda, con cifras reales

**[ICEX-CECO]** Aplicando el marco teórico que fundamentamos en el epígrafe 1.2: en 2022, la necesidad de financiación de las AAPP (4,81% PIB) más la adquisición neta de activos financieros (3,48% PIB) arroja un saldo neto a financiar del 8,29% PIB, que se cubrió íntegramente con variación neta de pasivos financieros por ese mismo importe; descontando los ajustes (2,63 puntos, que financian déficit pero no se computan como deuda bajo la metodología del Procedimiento de Déficit Excesivo), la **variación de la deuda PDE** fue de 5,66 puntos, con un stock final de **113,2% del PIB** a cierre de 2022. [Probable] Conviene contrastar este dato con la trayectoria histórica española previa a 2008, que el documento de referencia no ofrece pero que resulta imprescindible para calibrar la magnitud del cambio: España mantenía, en los años inmediatamente anteriores a la crisis financiera global, una de las ratios de deuda pública más bajas de toda la Unión Europea —en el entorno del 35-40% del PIB, muy por debajo del propio umbral de Maastricht—, de modo que el nivel actual, superior al 110%, representa, en apenas década y media, uno de los incrementos de endeudamiento más pronunciados de cualquier economía avanzada comparable, arrastrado sucesivamente por la crisis financiera de 2008-2013 y por la pandemia de 2020-2021.

### Los saldos por subsector: la Administración Central como amortiguador estructural

**[ICEX-CECO]** El desglose por subsector confirma, con cifras, el papel que ya anticipamos en el epígrafe 1.3: la **Administración Central** concentra, sistemáticamente, la mayor parte del ajuste —desde el -7,67% de 2020, en plena pandemia, hasta el -3,09% de 2022—, mientras que la **Administración Local** mantiene, desde 2012, superávits estructurales recurrentes (+0,31% en 2019, cayendo a -0,12% en 2022 por causas puntuales que desarrollamos a continuación), y la **Seguridad Social** un déficit moderado y decreciente. La **Administración Regional** pasó de un déficit contenido en 2021 (-0,05%) a un repunte notable en 2022 (-1,14%). [Probable] FEDEA, en su análisis de las cuentas autonómicas de 2022, atribuye este repunte no a un deterioro estructural genuino sino a **factores atípicos** —una inversión pública anormalmente baja ese año, y fuertes subvenciones a los intereses de la deuda canalizadas a través de los mecanismos estatales de financiación y liquidez que ya trabajamos en el Tema 21—, advirtiendo de que el déficit subyacente, corregido de esos factores puntuales, se habría reducido a casi la mitad — la misma cautela metodológica que ya hemos aplicado, en otros contextos de este proyecto, ante cualquier cifra anual aislada que pueda estar distorsionada por circunstancias no recurrentes.

### La deuda por subsector, y el caso valenciano como ilustración extrema de la desigualdad territorial

**[ICEX-CECO]** Del stock total de deuda PDE consolidada (113,2% PIB en 2022), la Administración Central concentra el 79,8%, la Regional el 23,9%, la Local el 1,7% y la Seguridad Social el 7,8%. [Seguro] Dentro del nivel regional, la dispersión entre Comunidades Autónomas es extraordinaria: Madrid, Canarias y los territorios forales apenas superan el 13-14% de deuda sobre su propio PIB regional, mientras que Murcia, Castilla-La Mancha y Cataluña rondan el 32-34%, y la **Comunidad Valenciana** alcanza el **44,4%** —con enorme diferencia, la más endeudada de toda España—. [Suposición] Este dato valenciano no es, según denuncia de forma reiterada y transversal la propia clase política de esa Comunidad desde hace más de una década, una simple cuestión de mala gestión autonómica: la Comunidad Valenciana sostiene que el sistema de financiación autonómica vigente desde 2009 —que ya trabajamos con todo detalle en el Tema 21— la sitúa de forma estructural entre las peor financiadas por habitante ajustado de toda España, obligándola a un recurso al endeudamiento sistemáticamente mayor que el de Comunidades con una financiación per cápita más generosa para sostener un nivel de servicios públicos comparable —la llamada **"infrafinanciación valenciana"**, un agravio que ha generado, en los últimos años, manifestaciones multitudinarias y una demanda transversal de reforma del modelo que sigue, como ya señalamos al hablar de la LOFCA y de la Ley 22/2009, pendiente de la revisión que la propia norma prevé mecanismos para acometer y que lleva más de quince años sin producirse. Es, en definitiva, la confirmación con datos de deuda —el indicador más duro y menos manipulable de todos los que hemos trabajado en este tema— de que la asimetría territorial española no es solo un debate político abstracto, sino una realidad financiera medible, con consecuencias muy concretas sobre la capacidad de cada Comunidad Autónoma de sostener su propio Estado del bienestar.

### Verificación final del Bloque 1

Repaso lo cubierto en las cuatro piezas de este bloque: la finalidad teórica de tratar a las AAPP como unidad económica, con el problema de medición de la producción no de mercado y el Informe Atkinson; el fundamento de por qué se aíslan las operaciones financieras, con el caso griego como ilustración dramática; el aparato teórico completo de los subsectores, con España como caso comparado único —"país regional" según la OCDE—; y esta situación actual, con el Informe de Envejecimiento de 2024 como pieza de sostenibilidad de mayor actualidad, el ajuste déficit-deuda con cifras reales, y el caso valenciano como cierre de la desigualdad territorial medida en deuda. No identifico ninguna vertiente adicional que el título de este bloque exija y que quede sin cubrir.



\clearpage
\section{Finanzas de las Administraciones Públicas: Estructura de ingresos}
\puntosclave{
    
}

## Introducción al Bloque 2

Cerrado el Bloque 1 —qué significa tratar a las AAPP como unidad económica, y cuál es su situación agregada y por subsector—, entramos en el primero de los dos lados de la cuenta: de dónde procede el dinero que sostiene toda esa actividad. Este bloque recorre dos preguntas distintas, y conviene fijar desde el principio que no son la misma pregunta formulada dos veces: el epígrafe 2.1 pregunta **de qué tipo** es cada ingreso —su estructura, su composición por naturaleza económica—; el epígrafe 2.2 preguntará **cuánto**, en conjunto, extrae el sector público de la economía —la presión fiscal—. Son, en cierto sentido, la misma relación que existe entre conocer la composición nutricional de un alimento y conocer cuántas calorías aporta en total: ambas preguntas son necesarias, y ninguna sustituye a la otra.

## 2.1. Estructura general de ingresos

### Por qué la estructura importa tanto como el nivel: una pregunta que el propio nombre del epígrafe no explicita

[Probable] Antes de recorrer las rúbricas concretas, merece la pena detenerse en algo que rara vez se explicita con la claridad que merece: **por qué debería importarnos de qué tipo es cada ingreso**, más allá de su magnitud total. La razón de fondo es que cada tipo de ingreso —impuesto directo, indirecto, cotización social— tiene **propiedades económicas radicalmente distintas** entre sí, algunas de las cuales ya trabajamos con detalle en otros temas de este proyecto y que aquí conviene recordar de forma sintética, precisamente porque su combinación es lo que da a un sistema fiscal su carácter propio. La sensibilidad al ciclo económico difiere: el Impuesto sobre Sociedades y el IRPF, ligados a beneficios y rentas, se contraen con fuerza en una recesión —el comportamiento anticíclico que ya trabajamos al hablar del servicio de la deuda en el Tema 22—, mientras que las cotizaciones sociales y el IVA resultan más estables. La movilidad internacional de la base imponible difiere —recuerda toda la teoría de Ramsey y de la elasticidad inversa que trabajamos en el Tema 13, y la propia tendencia europea a la "carrera a la baja" en el Impuesto sobre Sociedades que ya vimos en el Tema 17—. Y la **visibilidad** para el contribuyente difiere de forma decisiva —recuerda la teoría de la ilusión financiera de Puviani y la hipótesis de Mill que trabajamos con todo detalle en el Tema 13—. La estructura de ingresos no es, por tanto, un simple ejercicio de contabilidad descriptiva: es el mapa que revela cómo un país ha decidido repartir, entre instrumentos con propiedades económicas y políticas muy distintas, la financiación de su sector público.

### El argumento que da a este epígrafe su mayor profundidad: la estructura de ingresos como huella dactilar del modelo de Estado del bienestar

[Seguro] Existe, además, un argumento de mayor calado teórico que conecta directamente la pregunta "¿de qué tipo son los ingresos?" con una de las clasificaciones más influyentes de toda la ciencia política comparada del Estado del bienestar: la tipología de **Gøsta Esping-Andersen**, formulada en su obra ya clásica *The Three Worlds of Welfare Capitalism* (1990), que distingue tres grandes regímenes de bienestar según cómo organizan la protección social —**liberal** (Estados Unidos, Reino Unido, Canadá: prestaciones modestas, sujetas a comprobación de recursos, fuerte dependencia del mercado), **conservador-corporatista** (Alemania, Francia, Austria: protección social organizada en torno al seguro social vinculado al empleo, con prestaciones diferenciadas por estatus profesional) y **socialdemócrata** (los países nórdicos: cobertura universal, no vinculada al historial de cotización)—. [Probable] Lo que la literatura posterior a Esping-Andersen ha desarrollado con precisión, y que conecta de forma directa con este epígrafe, es que cada uno de estos tres modelos tiende a **financiarse de forma sistemáticamente distinta**: los regímenes corporatistas —los que la propia tradición alemana denomina, con el nombre de su fundador, **modelo bismarckiano**— se financian predominantemente mediante **cotizaciones sociales** vinculadas al empleo, reforzando la lógica de "seguro" y de contraprestación entre lo cotizado y lo percibido; los regímenes socialdemócratas y el modelo británico —la tradición que toma su nombre de William Beveridge, arquitecto del Estado del bienestar británico de posguerra— se financian, en cambio, predominantemente mediante **impuestos generales**, coherente con su vocación de cobertura universal desvinculada de la trayectoria laboral individual de cada beneficiario.

[Seguro] Aplicado a España, este marco teórico no es una abstracción sin consecuencias: **confirma, con toda precisión, el carácter bismarckiano-corporatista** de nuestro sistema de protección social que ya intuíamos, sin nombrarlo así, al trabajar la Seguridad Social en el Tema 22 —recuerda el propio nombre de la institución, "Seguridad **Social**", y su financiación mediante cotizaciones vinculadas al empleo—. La propia estructura de ingresos que recorremos a continuación es, en este sentido, la confirmación estadística exacta de esa filiación teórica: España comparte, con Alemania y Francia, un peso de las cotizaciones sociales muy superior al de los países anglosajones o nórdicos, precisamente porque comparte con ellos la misma tradición institucional de origen bismarckiano en la construcción de su Estado del bienestar.

### Las rúbricas concretas: la jerarquía que ya conocíamos, ahora con su fundamento completo

**[ICEX-CECO]** Con este marco ya fijado, las cifras de 2022 cobran su sentido completo. Las **cotizaciones sociales** —13,6% del PIB (efectivas de empleadores, 9,6%; efectivas de hogares, 3,4%; imputadas, 0,5%)— son la primera rúbrica de ingresos de las AAPP españolas, confirmando con la máxima claridad posible el carácter bismarckiano que acabamos de fundamentar. Los **impuestos directos** —12,4% del PIB, de los que los impuestos sobre la renta representan 12,1 puntos, repartidos aproximadamente en un 80% entre personas físicas y un 20% entre sociedades— ocupan el segundo lugar, con un peso muy próximo al de las propias cotizaciones —una cercanía que, conviene señalarlo, no es universal entre los países de tradición bismarckiana, y que revela que España, aun compartiendo la filiación corporatista de fondo, ha desarrollado también un pilar de imposición directa comparativamente robusto—. Los **impuestos indirectos** —12,1% del PIB, de los que el IVA aporta 7,1 puntos y el resto (Impuestos Especiales, ITPAJD, y demás figuras ya trabajadas con detalle en los Temas 17, 19 y 20) los 3,2 puntos restantes, más 1,8 puntos adicionales de "otros impuestos sobre la producción" (IBI, IAE, vehículos de tracción mecánica)— completan el tercer gran pilar. Los **otros ingresos corrientes** (rentas de la propiedad, ventas de bienes y servicios, transferencias diversas) aportan un 3,8% adicional, y los **impuestos sobre el capital** (sucesiones y donaciones, plusvalía municipal, contribuciones especiales) apenas un 0,4% —el nivel que la propia serie histórica española considera "normal", sin apenas variación relevante de un ejercicio a otro—.

### Verificación de cierre

[Probable] Este epígrafe ha cubierto: por qué la estructura de ingresos merece atención propia, distinta de su nivel agregado, dadas las propiedades económicas diferenciadas de cada tipo de ingreso; el marco teórico de Esping-Andersen y la distinción bismarckiano/beveridgiano aplicada, con precisión, a la propia filiación institucional española ya trabajada en los Temas 21 y 22; y el desglose completo de las cinco grandes rúbricas de ingreso en contabilidad nacional con sus cifras de 2022. Queda, para el epígrafe siguiente, la pregunta complementaria del nivel agregado —la presión fiscal—, que este epígrafe ha dejado deliberadamente para después, tal y como acordamos al fijar la estructura del bloque.


\subsection{Presión fiscal: qué es, cómo se mide, comparativa internacional, historia y desglose}

### Qué es, y por qué no es lo mismo que "impuestos pagados"

[Seguro] La presión o carga fiscal global se define como los impuestos y cotizaciones sociales recaudados por las AAPP y por las instituciones europeas, más los impuestos que los residentes de un país pagan al resto del mundo, expresados en proporción del PIB. [Probable] Conviene precisar, porque el epígrafe anterior ya nos dio las herramientas para hacerlo, que esta magnitud **no es la simple suma** de las rúbricas de ingreso que acabamos de recorrer: incluye, además, el ajuste por **recaudación incierta** —la corrección técnica que aplica a los impuestos y cotizaciones devengados pero cuyo cobro efectivo resulta dudoso, ya mencionada en el epígrafe sobre transferencias de capital del Bloque 1 original— y la parte de recaudación que se transfiere directamente a las instituciones de la Unión Europea, que forma parte de la presión fiscal que soporta el contribuyente aunque nunca llegue a integrarse en el presupuesto puramente nacional. Dos organismos internacionales calculan esta magnitud con metodologías no plenamente coincidentes pero de resultados muy próximos: la **Comisión Europea**, conforme al Reglamento Delegado (UE) nº 877/2013, y la **OCDE**, en su serie *Revenue Statistics* — esta última, conviene señalarlo, una de las series estadísticas más longevas de todo el organismo, publicada de forma ininterrumpida desde **1965**, lo que la convierte en la referencia comparativa internacional más consolidada y de mayor recorrido histórico de cuantas hemos manejado en todo este proyecto.

### La comparación internacional: España, en el pelotón intermedio de Europa

**[ICEX-CECO]** Según la base AMECO de la Comisión Europea, la presión fiscal española en 2022 se situó en el **38,23%** del PIB, frente al **40,42%** de la Unión Europea en su conjunto y al 40,77% de la UE-15 — España, por tanto, con una fiscalidad relativamente baja dentro del contexto europeo, aunque no entre los países de menor presión fiscal del continente. Los datos de la OCDE para 2021 sitúan a Dinamarca (46,9%), Francia (45,1%) y Suecia (42,6%) a la cabeza, seguidos por la media de la UE (40,8%) y Alemania (39,5%); España (38,4%) se sitúa justo por debajo de esta última, seguida por Nueva Zelanda, Reino Unido y Australia entre los países anglosajones, y cerrando el grupo Estados Unidos (26,6%), Chile (22,2%), Irlanda (21,1% —caso "atípico" fiscal por su renta per cápita anormalmente elevada, ya señalado en el documento original—) y México (16,7%).

### La historia reciente: por qué España ha subido cuando el resto de Europa bajaba

**[ICEX-CECO]** El dato de mayor interés analítico de todo este epígrafe es el de la **trayectoria**, no solo el del nivel: España experimentó un incremento de **3,4 puntos** de presión fiscal entre 2019 y 2022, mientras la UE en su conjunto **caía** 0,12 puntos y la UE-15 caía 0,91 —una divergencia de dirección, no solo de magnitud, que merece explicación—. El propio Programa de Estabilidad 2023-2026 preveía superar el umbral del 40% del PIB en 2024. [Probable] Las razones que el propio documento de referencia apunta —que la recaudación tributaria creció, en los años de recuperación pospandemia, más deprisa que la demanda interior nominal (en 2022, +14,4% de recaudación frente a +12,0% de demanda), el afloramiento de bases imponibles durante la Covid, y la caída del desempleo estructural (NAWRU del 11,0% en 2023 frente al 12,9% de 2019)— son correctas pero incompletas, porque dejan fuera un mecanismo concreto, muy debatido en la política fiscal española reciente, que conviene desarrollar con precisión: la **no deflactación del IRPF**, o lo que la propia literatura técnica denomina **"progresividad en frío"** (el equivalente español del concepto internacional de *bracket creep* o *fiscal drag*).

[Seguro] El mecanismo es sencillo de explicar y ya lo hemos trabajado, en esencia, al hablar del cambio de tramo de un contribuyente cuyo salario sube exactamente al ritmo de la inflación: si los tramos del IRPF permanecen fijos en términos nominales mientras los salarios se revisan para compensar la subida de precios, una parte creciente de esa renta —que en términos **reales** no ha aumentado en absoluto— pasa a tributar a un tipo marginal superior, generando una subida de la recaudación —y, por tanto, de la presión fiscal medida en términos agregados— **sin que exista ninguna decisión legislativa expresa de subir impuestos**. [Seguro] España solo ha "deflactado" la tarifa estatal del IRPF en dos ocasiones en los últimos veinte años —en 2008, bajo el Gobierno de José Luis Rodríguez Zapatero, y en 2015, bajo el de Mariano Rajoy, esta última en el marco de una reforma fiscal más amplia—, y ambas coincidieron, no por casualidad según señalan los propios analistas del episodio, con la proximidad de sendas convocatorias electorales generales. [Seguro] Durante el repunte inflacionario de 2021-2023 —con una inflación interanual que llegó a superar el 10% en 2022, el nivel más alto desde 1984—, la tarifa **estatal** del IRPF no se deflactó en ningún ejercicio, generando un intenso debate político: la oposición sostuvo que la falta de ajuste equivalía a una "subida de impuestos encubierta"; el Gobierno, por boca del propio director general de la Agencia Tributaria, Jesús Gascón, defendió que deflactar beneficiaría desproporcionadamente a las rentas más altas —los contribuyentes de menor renta ya están exentos del impuesto y, por tanto, no se benefician de ningún ajuste de tramos— y que reduciría los ingresos públicos en un momento de necesidad presupuestaria, además de poder alimentar aún más la propia inflación. [Probable] Conviene, con la misma ecuanimidad que hemos aplicado a cada debate político de este proyecto, señalar que ambos argumentos tienen fundamento técnico real, y que se trata de uno de esos debates de economía pública donde la evidencia no zanja, por sí sola, cuál es la decisión "correcta" —depende, en última instancia, de qué prioridad distributiva y qué prioridad de disciplina fiscal se pondere más—. [Probable] Sí conviene señalar, como dato comparado dentro de la propia España, que el **País Vasco** fue pionero en aplicar la deflactación en su propio tramo foral —un 5,5% acumulado durante 2022, entre la propia Comunidad Autónoma y las Diputaciones Forales—, y que la Comunidad de Madrid ha sido, entre las de régimen común, la que con mayor regularidad ha deflactado su propio tramo autonómico en los últimos ejercicios — un ejemplo más, aplicado ahora al terreno fiscal más sensible de todos, de la capacidad normativa diferenciada entre Comunidades Autónomas que ya trabajamos con detalle en el Tema 21.

### El desglose detallado: por tipo de impuesto, por nivel de gobierno, y por función económica

**[ICEX-CECO]** Por tipo de gravamen (OCDE 2022, datos 2021), España distribuye su presión fiscal así: Renta y Beneficios 29,9%, Propiedad 7,1%, Seguridad Social 35,6%, Bienes y Servicios 27,4% — confirmando, una vez más, el peso dominante de las cotizaciones que ya fundamentamos con Esping-Andersen en el epígrafe anterior, y un patrón muy próximo al de Francia y Alemania, frente al escaso peso de las cotizaciones en Dinamarca, Australia o Nueva Zelanda —países que financian su protección social predominantemente mediante impuestos generales, coherente con su filiación beveridgiana o directamente liberal—.

**[ICEX-CECO]** Con mayor detalle interno (año 2020): dentro de la renta, las personas físicas aportan el 81,6% frente al 18,4% de las sociedades; dentro de las cotizaciones, los empleadores aportan el 72,8%, los empleados el 14,4%, los autónomos el 7,5% y los desempleados el 5,2%; dentro de la propiedad, los impuestos recurrentes sobre bienes inmuebles suponen el 47,7%, las transacciones financieras y de capital (ITP, AJD) el 27,3%, sucesiones y donaciones el 9,0%; dentro de bienes y servicios, el tipo IVA aporta el 64,1%, los impuestos sobre determinados bienes y servicios (Impuestos Especiales, IEDMT, tasas del juego) el 27,4%.

**[ICEX-CECO]** Por nivel de gobierno (OCDE 2022, datos 2020): Administración Central 37,6%, Administración Regional 16,7%, Administración Local 8,9%, Fondos de la Seguridad Social 36,8%. [Probable] Conviene recordar aquí, sin redesarrollarla, la precisión metodológica que ya fijamos en el documento de referencia original: la baja participación aparente de la Administración Regional en bienes y servicios (solo 4,9%) se debe al criterio OCDE de atribuir el ingreso al gobierno que "ejerce la autoridad de imponer el tributo" y "tiene discrecionalidad para variar el tipo" —IVA e Impuestos Especiales se atribuyen al Estado por ser este quien fija los tipos, con la excepción foral ya conocida—, no a que las Comunidades Autónomas no reciban, de hecho, una parte sustancial de esa recaudación por la vía de las transferencias del sistema de financiación que ya trabajamos en el Tema 21.

**[ICEX-CECO]** Por función económica (Comisión Europea 2023a, datos 2021): Consumo 9,8% del PIB, Trabajo 19,5% (empleadores 9,9, empleados 7,8, desempleados 1,9), Capital 9,0% (sociedades 2,7, familias 1,0, autónomos 1,8, stock de capital 3,2), total 38,4%. Los impuestos sobre el trabajo suponen así el **50,8%** de toda la presión fiscal española —prácticamente idéntico al 51,4% de la media de la UE—, con un tipo impositivo implícito sobre el trabajo del 36,5% en España frente al 37,8% de la UE.

### La cuña fiscal del trabajo: la magnitud que sintetiza todo lo anterior en una sola cifra

**[ICEX-CECO]** Cierro con el concepto que Fernández, Ponz y Taguas (1995) formularon para la Dirección General de Presupuestos española: la **cuña fiscal**, la diferencia entre el coste salarial que paga el empleador y el salario neto que percibe el trabajador, expresada como porcentaje del coste salarial bruto —fórmula Δ = 100 − [(1−r_l)(1−cs_T)/(1+cs_E)]×100, donde r_l es el tipo medio sobre rentas del trabajo, cs_T la cotización del trabajador y cs_E la del empresario—. Según la Comisión Europea (2023a), la cuña fiscal española en 2021 para un soltero sin hijos con salario medio fue del **27,9%**, frente al **31,9%** de la media de la UE. [Probable] Conviene señalar aquí una referencia internacional adicional que el documento original no cita, pero que constituye la fuente comparativa de mayor autoridad para este mismo indicador: la propia OCDE publica anualmente, desde hace décadas, su serie ***Taxing Wages***, que calcula la cuña fiscal de forma homogénea para todos sus países miembros y para distintos perfiles de hogar —soltero sin hijos, matrimonio con hijos, distintos niveles salariales—, siendo la referencia que cualquier análisis riguroso de competitividad fiscal sobre las rentas del trabajo debería consultar junto a los datos ya citados de la Comisión Europea.

[Seguro] El refinamiento del Modelo EREMS de Boscá, Doménech, Ferri y Méndez (2020), incorporando el poder adquisitivo mediante la imposición indirecta sobre el consumo (con un tipo medio calibrado τ_c = 0,15), eleva la cuña fiscal modificada de España hasta el **37,3%** — una diferencia de casi diez puntos respecto a la cuña fiscal simple, que revela cuánto se subestima el verdadero coste fiscal total del trabajo cuando se ignora la carga adicional que recae sobre el propio consumo del salario ya percibido. [Probable] Y la precisión final, ya apuntada en el documento de referencia: dado que lo relevante para la decisión de oferta de trabajo —recuerda, de nuevo, a Prescott y su explicación de por qué los europeos trabajan menos horas que los estadounidenses, ya trabajada en el documento original— son los **tipos marginales** y no los tipos medios, Boscá et al. (2022) estiman que, aplicando esta metodología más exigente, la cuña fiscal modificada española **superaría el 45%** — una cifra que, de confirmarse con el rigor metodológico que merece, sitúa el verdadero coste fiscal marginal del trabajo español en una magnitud muy superior a la que cualquiera de las comparaciones internacionales más simples, basadas en tipos medios, permite apreciar a primera vista.



Repaso lo cubierto en los dos epígrafes de este bloque: la estructura de ingresos con su fundamento en Esping-Andersen y la filiación bismarckiana española; y esta presión fiscal, con su definición y metodología de medición, su comparación internacional completa, su trayectoria reciente explicada con el mecanismo concreto de la progresividad en frío y el debate político real que genera, su desglose exhaustivo por tipo de impuesto, por nivel de gobierno y por función económica, y su síntesis final en la cuña fiscal del trabajo con sus distintos refinamientos metodológicos. No queda, dentro de la estructura de dos epígrafes que fijamos para este bloque, ninguna vertiente pendiente —recuerda que la desagregación por subsectores, que originalmente pensamos incluir aquí, quedó ya resuelta al fusionarse dentro del propio epígrafe 2.1, dado que la estructura de ingresos por naturaleza económica ya incorpora, de forma natural, la distinción entre lo que recauda cada nivel de gobierno.


\clearpage
\section{Finanzas de las Administraciones Públicas: Estructura de gastos}
\puntosclave{
    
}

Cerramos el tema, y con él toda esta serie de tres, con el último gran pilar: la estructura del gasto de las AAPP y su clasificación funcional internacional, la COFOG. Seguimos, en este bloque, exactamente el mismo patrón metodológico que ya aplicamos en el Tema 22 a la clasificación orgánica/económica/funcional presupuestaria: primero el fundamento teórico de la COFOG, sin cifras todavía (3.1); después la estructura general del gasto por naturaleza económica en contabilidad nacional, simétrica a la de ingresos que ya trabajamos (3.2); y, por último, la aplicación de la propia COFOG con cifras reales, agregada y por subsector (3.3).

\subsection{El aparato teórico de la COFOG}
El eslabón que faltaba en la genealogía que ya empezamos a construir en el Tema 22

[Seguro] Recordarás que, al fundamentar la clasificación funcional presupuestaria española en el Tema 22, situamos su origen en el manual de la ONU de 1958 sobre clasificación económica y funcional de las transacciones del Estado. La COFOG —Classification of the Functions of Government— es, en un sentido muy preciso que conviene explicitar ahora, el descendiente directo y hoy vigente de aquel manual de 1958: no una clasificación distinta con un origen independiente, sino la evolución formalizada de la misma idea, adaptada a las exigencias de comparabilidad internacional que el propio Sistema de Cuentas Nacionales fue exigiendo con el paso de las décadas.

Genealogía precisa: de 1980 a 1999, y una revisión que lleva veinticinco años pendiente

[Seguro] La COFOG se publicó, en su primera versión, en 1980; fue revisada en profundidad en 1999 —la versión conocida como COFOG-1999, aprobada por la Comisión de Estadística de las Naciones Unidas en su trigésima sesión, celebrada en marzo de ese mismo año, como parte de las cuatro clasificaciones básicas de gasto por finalidad que el SCN 1993 estableció—. [Seguro] Su desarrollo técnico fue obra conjunta de la OCDE, trabajando en estrecha colaboración con Eurostat, y de la propia División de Estadística de las Naciones Unidas —un reparto de responsabilidades dentro de la familia completa de cuatro clasificaciones de gasto por finalidad que el SCN 1993 estableció simultáneamente: la propia COFOG para el gasto de las Administraciones Públicas, la COPNI (Clasificación de las Finalidades de las Instituciones sin Fines de Lucro al Servicio de los Hogares), la COICOP (Clasificación del Consumo Individual por Finalidades, para el gasto de los hogares) y la COPP (Clasificación de las Erogaciones de los Productores por Finalidades) —, cuatro piezas hermanas de un mismo proyecto de estandarización internacional del gasto por propósito, de las que la COFOG es, con diferencia, la más citada y la más utilizada en el análisis de finanzas públicas comparadas—. [Seguro] La COFOG se estructura en tres niveles de detalle: divisiones (código de dos dígitos, 01 a 10, describiendo los grandes objetivos de la actuación pública), grupos (código de tres dígitos, describiendo los servicios concretos mediante los que se persigue cada objetivo) y clases (código de cuatro dígitos, el nivel de máximo detalle) — recordarás, del Tema 21 original, que España no implementa habitualmente este tercer nivel de clases en sus propias estadísticas.

[Seguro] Conviene cerrar esta genealogía con un dato de extraordinaria actualidad que ningún manual anterior a esta misma redacción podría haber incluido: pese a que el panorama de las políticas públicas y de los objetivos socioeconómicos de los Estados ha cambiado de forma sustancial en el último cuarto de siglo —piensa, sin ir más lejos, en toda la nueva fiscalidad ambiental y en las políticas de transición energética que hemos trabajado en el Tema 20, ausentes por completo del horizonte de 1999—, la COFOG no ha sido revisada desde entonces. La propia Comisión de Estadística de las Naciones Unidas, en su quincuagésimo tercera sesión, celebrada en 2022, decidió formalmente iniciar su revisión, encomendándola a un equipo de trabajo específico —el TT-COFOG— constituido bajo el Comité de Expertos de las Naciones Unidas en Clasificaciones Estadísticas Internacionales, con el mandato expreso de desarrollar una versión actualizada de la clasificación. [Probable] Es, por tanto, un proceso de reforma genuinamente en curso en este mismo momento, y conviene que sepas que cualquier fecha de examen posterior a la finalización de este trabajo debería consultar si esa revisión ya ha visto la luz, y con qué contenido, antes de dar por definitiva la estructura de diez divisiones que trabajaremos con cifras en el epígrafe 3.3.

Por qué es la clasificación de referencia internacional: el mismo argumento de comparabilidad, ahora aplicado al gasto por finalidad

[Probable] La razón por la que la COFOG, y no cualquiera de las múltiples clasificaciones funcionales nacionales que cada país desarrolla para su propio uso doméstico —recuerda, del Tema 22, las 27 políticas de gasto españolas agrupadas en 5 áreas—, se ha convertido en la referencia obligada de cualquier análisis comparado de finanzas públicas, es exactamente el mismo argumento de fondo que ya fundamentamos, con Kevin O'Connor y el propio Manual de Estadísticas de Finanzas Públicas del FMI, al hablar de por qué la consolidación intersectorial es indispensable para la comparabilidad entre países: las clasificaciones funcionales domésticas de cada país no son comparables entre sí, porque cada una responde a la estructura administrativa, a las prioridades políticas y a la tradición estadística propias de ese país concreto —ya vimos, precisamente, que la propia área de "Servicios Públicos Básicos" española tiene un contenido sustancialmente distinto según se aplique al Estado o a las Entidades Locales—. La COFOG, al ser una clasificación única, idéntica y obligatoria para todos los países que remiten estadísticas a Eurostat, a la OCDE, al FMI o a Naciones Unidas, es la única herramienta que permite comparar, con la confianza de estar midiendo exactamente la misma cosa, cuánto gasta Francia en sanidad frente a cuánto gasta España, o cómo ha evolucionado el peso de la defensa en el gasto público de cada país europeo a lo largo de las últimas décadas.

Ventajas e inconvenientes: la misma tensión de fondo que ya conocemos, con su propia variante

La ventaja de la COFOG es la comparabilidad internacional perfecta. Su inconveniente principal es el reverso de esa virtud: precisamente por estar diseñada para funcionar de forma idéntica en decenas de países con estructuras institucionales y prioridades políticas muy distintas entre sí, la COFOG puede resultar, en la práctica doméstica de cualquier país concreto, menos útil para la gestión interna que una clasificación nacional a medida —la misma tensión entre estandarización internacional y ajuste institucional doméstico que ya identificamos, en otro contexto, al hablar de la clasificación por programas española del Tema 22—. [Suposición] Un segundo inconveniente, ya apuntado y que conviene retener como cierre de este epígrafe, es que la COFOG —como cualquier clasificación funcional, según ya trabajamos con el problema de las fronteras arbitrarias en el Tema 22— no está exenta de dificultades de encaje para políticas genuinamente transversales: la propia transición ecológica, que hoy atraviesa simultáneamente las divisiones de protección del medio ambiente, asuntos económicos, y vivienda y servicios comunitarios, es difícil de rastrear como magnitud única y coherente dentro de un esquema diseñado, en su formulación vigente, hace más de un cuarto de siglo — precisamente la carencia que la revisión en curso de 2022 pretende, entre otros objetivos, subsanar. Y un tercer inconveniente, que conecta directamente con el Informe Atkinson ya trabajado en el Bloque 1 de este mismo tema: la COFOG clasifica el gasto por finalidad, no por resultado ni por eficiencia —dos países pueden gastar exactamente el mismo porcentaje de su PIB en la división "Salud", clasificados de forma idéntica bajo COFOG, y obtener resultados sanitarios radicalmente distintos, sin que la propia clasificación tenga ninguna capacidad de capturar esa diferencia—, la misma limitación de fondo que ya identificamos al hablar de la convención "output = input" y de la dificultad estructural de medir la verdadera productividad del sector público.


\subsection{Estructura general de gastos}
La simetría con el Bloque 2, aplicada ahora al otro lado de la cuenta

[Seguro] Ya lo dejamos anunciado al resolver tu pregunta sobre la simetría de bloques: esta es la misma clasificación por naturaleza económica que ya trabajamos para los ingresos, aplicada ahora al gasto —transferencias sociales, remuneración de asalariados, consumos intermedios, intereses, subvenciones, formación bruta de capital, transferencias de capital—, no la COFOG que acabamos de fundamentar en el epígrafe anterior. Es, dentro de la secuencia de cuentas del SCN que ya conocemos por Richard Stone, la misma pregunta que hacíamos del lado del ingreso —¿de qué naturaleza económica es esta partida?—, formulada ahora del lado del gasto.

Transferencias sociales: la rúbrica dominante, y la teoría económica que explica por qué crece más deprisa que casi todo lo demás

[ICEX-CECO] Las transferencias sociales son, con diferencia, la rúbrica más importante: 20,1% del PIB en 2022, de los que 17,2 puntos corresponden a prestaciones distintas de las de especie (pensiones contributivas y no contributivas de la SS, 12,9 puntos; Clases Pasivas del Estado, 2,0; desempleo no vinculado a la Covid, 1,6; incapacidad laboral transitoria, 1,0 —con el añadido excepcional de 4,4 puntos por ERTES en 2021, ya trabajados con detalle en el Tema 22, que cayeron un 90% en 2022—) y 2,9 puntos a transferencias sociales en especie (educación concertada, convenios sanitarios privados, farmacia no hospitalaria, mutualismo administrativo de funcionarios —recuerda MUFACE, ISFAS y MUGEJU, ya trabajados en el Tema 22—). [ICEX-CECO] Su crecimiento entre 2008 y 2022 —+67,9% frente a un +30,3% de la remuneración de asalariados en el mismo periodo— es, según reconoce el propio documento de referencia, el principal responsable del incremento del gasto corriente español, y ya conectamos esta cifra, en el Bloque 1, con el Informe de Envejecimiento de 2024 y la propia reforma de pensiones de 2021-2023, que la propia Comisión Europea señala como la que más agravará esta misma tendencia de cara a las próximas décadas. No repito aquí ese desarrollo, ya construido con todo el detalle que merece; me limito a señalar la coherencia completa entre ambas piezas del tema.

Remuneración de asalariados y consumos intermedios: la teoría que explica por qué ni siquiera una gestión pública perfecta evitaría su encarecimiento estructural

[ICEX-CECO] La remuneración de asalariados (11,6% PIB) y los consumos intermedios (5,9% PIB, con el ejemplo ya señalado de los medicamentos hospitalarios, más caros que en la farmacia ordinaria por incluir tratamientos oncológicos y anestesias de alta complejidad) son la segunda y la tercera rúbrica del gasto corriente. [Seguro] Existe una teoría económica precisa, y muy citada en la literatura de finanzas públicas comparadas, que explica por qué estas dos partidas concretas —las que sostienen, sobre todo, la sanidad y la educación públicas— tienden a crecer de forma estructural por encima de la inflación general, con independencia de cuán eficiente sea su gestión: la llamada enfermedad de los costes de Baumol (Baumol's cost disease), formulada por William Baumol y William Bowen en los años sesenta a partir de su estudio sobre la economía de las artes escénicas. [Seguro] Su intuición original —que un cuarteto de cuerda necesita hoy exactamente el mismo número de músicos y el mismo tiempo para interpretar una pieza que hace un siglo, sin ninguna ganancia de productividad posible— se traslada, con notable fuerza explicativa, a los servicios públicos intensivos en trabajo humano: una enfermera no puede cambiar un vendaje más deprisa hoy que hace cincuenta años, ni un profesor puede corregir un examen con mayor rapidez, por mucho que la tecnología haya transformado, en cambio, la productividad de sectores manufactureros comparables. [Probable] Como esos trabajadores compiten, en el mercado laboral, con sectores donde sí existe ganancia de productividad —y, por tanto, capacidad de pagar salarios crecientes—, los salarios de la sanidad y la educación deben subir al mismo ritmo, sin que exista ninguna ganancia de productividad equivalente que lo justifique en términos puramente técnicos, generando un encarecimiento estructural y persistente de estos servicios que ninguna reforma de gestión, por acertada que sea, puede eliminar por completo —solo mitigar—. [Seguro] La magnitud de lo que está en juego es considerable: la propia literatura que ha aplicado esta teoría al gasto público señala que sanidad y educación representaban ya, en 2006, el 11,1% del PIB conjunto de los países de la OCDE, más de una cuarta parte de todo el gasto público de esos países. [Probable] Conviene, con la misma honestidad analítica de siempre, señalar que esta teoría no está exenta de matización crítica: Ferris y West, en un trabajo publicado en Public Choice en 1996, cualificaron la aplicación mecánica de Baumol al crecimiento del sector público en su conjunto, señalando que otros factores —de naturaleza más política que puramente tecnológica, como la propia demanda ciudadana creciente de servicios públicos de calidad— explican también una parte sustancial de esta misma tendencia, sin que la enfermedad de los costes sea, por sí sola, una explicación completa y suficiente.

Intereses: el mismo comportamiento anticíclico, ya conocido

[ICEX-CECO] Los intereses de la deuda (2,4% PIB en 2022) muestran el comportamiento anticíclico ya trabajado en el Tema 22 —bajan en expansión, suben en recesión—, sin que el comportamiento igualmente procíclico de los tipos de interés de mercado logre, según señala el documento original, compensar por completo ese patrón.

Subvenciones: del apoyo genérico a la respuesta específica ante la crisis energética, con su marco europeo propio

[ICEX-CECO] Las subvenciones (2% PIB en 2022) reflejan, sobre todo, las ayudas y bonificaciones a productos energéticos y de transporte derivadas de la guerra de Ucrania, sin ningún patrón de tendencia sostenida al alza en el resto de la serie histórica. [Probable] Conviene añadir aquí el marco jurídico europeo que hizo posible este tipo de subvención masiva sin infringir la disciplina general de ayudas de Estado que ya trabajamos, en otro contexto, al hablar del criterio del inversor privado en el anexo del Tema 22: la Comisión Europea aprobó, en marzo de 2022, un Marco Temporal de Crisis para ayudas de Estado, que flexibilizó de forma excepcional y acotada en el tiempo los límites ordinarios que el art. 107 TFUE impone a las subvenciones públicas a empresas y consumidores, precisamente para permitir a los Estados miembros amortiguar el impacto del encarecimiento energético derivado de la invasión rusa sin necesidad de recurrir, caso por caso, al procedimiento ordinario de notificación y autorización previa —el mismo tipo de flexibilización excepcional y temporal que ya vimos operar, con otro instrumento, en la cláusula de escape de las reglas fiscales que trabajamos en el Tema 21.

Formación bruta de capital: el gasto que paga el precio político de ser el más fácil de recortar

[ICEX-CECO] La inversión pública (2,8% PIB en 2022) muestra una trayectoria descendente de largo plazo —del entorno del 5% del PIB a mediados de los noventa, en pleno auge de los Fondos Estructurales y de Cohesión europeos, a 4,4 puntos en 2004-2008, 3,8 en 2009-2013 y apenas 2,2 en 2014-2018—, con una comparación internacional desfavorable: Eurozona 3,0%, UE y EEUU 3,2%, Japón 4,2%. [Probable] El propio documento de referencia señala, sin desarrollarlo, que se trata de "gasto discrecional" que "sufre recortes" en los ajustes urgentes — una observación que la literatura internacional sobre consolidación fiscal ha documentado con notable consistencia bajo el nombre de sesgo hacia la inversión en los procesos de ajuste presupuestario: cuando un Gobierno necesita recortar gasto con rapidez, la inversión pública resulta, políticamente, el objetivo de menor coste inmediato —a diferencia de recortar pensiones, salarios públicos o prestaciones sociales, que golpean de forma visible e inmediata a colectivos concretos de votantes, posponer o cancelar un proyecto de infraestructura no genera, en el corto plazo electoral, ningún perjuicio identificable para nadie en particular, aunque sus consecuencias sobre el crecimiento potencial de largo plazo puedan ser, con el tiempo, sustanciales—. Es, en el fondo, la misma lógica de descuento temporal excesivo del beneficio político inmediato frente al coste diferido que ya trabajamos, con Kydland y Prescott, al hablar de la inconsistencia temporal en el Tema 21: la inversión pública, precisamente por no tener un beneficiario políticamente organizado y visible que la defienda en el corto plazo, es la primera víctima de cualquier ajuste fiscal urgente, con independencia de cuál sea, en términos estrictamente económicos, su verdadero valor de largo plazo para el crecimiento del país.

Transferencias de capital: el catálogo ya conocido, sin necesidad de redesarrollo

[ICEX-CECO] Las transferencias de capital (1,1% PIB en 2022) agrupan las ayudas a la inversión de terceros y otras transferencias de redistribución patrimonial —ayudas al sector financiero vía FROB-SAREB, ayuda al desarrollo, préstamos declarados incobrables y avales ejecutados, sentencias judiciales y devoluciones tributarias—, sin tendencia sostenida al alza, ya trabajadas con suficiente detalle en el propio documento de referencia sin necesidad de ampliación adicional en este punto.

Verificación final del epígrafe

Este epígrafe ha cubierto la totalidad de las rúbricas económicas del gasto, con dos aportaciones teóricas genuinamente nuevas que el documento de referencia no ofrecía: la enfermedad de los costes de Baumol como explicación estructural del crecimiento de remuneración de asalariados y consumos intermedios, con su propia matización crítica; y el sesgo hacia la inversión como explicación de política económica del comportamiento cíclico de la formación bruta de capital. No identifico ninguna vertiente adicional pendiente dentro del alcance que este epígrafe se propuso.


\subsection{COFOG aplicada: cifras y porcentajes de PIB por función}
La tabla completa

[ICEX-CECO]

División COFOG (% PIB)	2021	2022	2023(p)	2024(p)
01. Servicios públicos generales	5,9	5,8	5,7	5,7
02. Defensa	1,0	1,1	1,2	1,2
03. Orden público y seguridad	2,0	1,9	1,9	1,9
04. Asuntos económicos	6,3	5,5	5,5	4,7
05. Protección del medio ambiente	1,0	1,0	1,0	1,0
06. Vivienda y servicios comunitarios	0,5	0,5	0,5	0,5
07. Sanidad	7,3	6,9	6,9	6,9
08. Ocio, cultura y religión	1,2	1,2	1,2	1,1
09. Educación	4,6	4,5	4,5	4,5
10. Protección social	20,6	18,9	19,0	19,0
TOTAL	50,4	47,4	47,3	46,3

El gasto social en sentido amplio (Protección Social + Sanidad + Educación + Ocio/Cultura + Vivienda) alcanza el 33% del PIB, el 69,6% del gasto total —confirmando, con la clasificación internacional de referencia y no ya con la clasificación por políticas puramente española del Tema 22, la misma centralidad del gasto social que ya documentamos entonces.

Por subsectores: la aplicación directa de todo lo que ya sabemos

[Probable] No necesitamos cifras adicionales para saber, con razonable seguridad, cómo se reparte cada división entre los cuatro subsectores: Sanidad y Educación, dado el grado de descentralización superior al 90% que ya cuantificamos en el Bloque 1, corresponden abrumadoramente a la Administración Regional; Protección Social, por la propia naturaleza funcional y no territorial de la Seguridad Social que fundamentamos en aquel mismo bloque, corresponde casi en su totalidad a los Fondos de la Seguridad Social; Defensa, Servicios Públicos Generales —con la Deuda Pública incluida— y Orden Público son, por su propia naturaleza de soberanía indivisible, competencia casi exclusiva de la Administración Central; y Vivienda y Servicios Comunitarios recae, en su mayor parte, sobre la Administración Local, coherente con las competencias urbanísticas municipales ya trabajadas en el Tema 21.

El declive de "Asuntos económicos": el cierre exacto del círculo que abrimos con el PRTR en el Tema 22

[Probable] La división que muestra la caída más pronunciada de toda la tabla es "Asuntos Económicos" —de 6,3 a 4,7 puntos de PIB entre 2021 y 2024, un descenso de 1,6 puntos—, y esta caída no es casual ni ajena a todo lo que ya construimos en el tema anterior: recuerda que el 71,7% del PRTR y REACT-EU se destinaba precisamente a Actuaciones de Carácter Económico —I+D+i, infraestructuras, industria y energía—, por el mandato regulatorio europeo del 37%/20% que fundamentamos con todo detalle. A medida que los fondos del Mecanismo de Recuperación y Resiliencia se van ejecutando y aproximándose a su fecha límite —recuerda, el 70% de las ayudas no reembolsables debía estar jurídicamente comprometido antes de finales de 2022, y el conjunto del Mecanismo se agota en el horizonte de mediados de la década—, el gasto excepcional que esos fondos inyectaban específicamente en esta división se va desvaneciendo, devolviendo la serie histórica hacia niveles más próximos a los que tenía antes de 2021. Es, literalmente, el cierre del círculo que abrimos hace dos temas: la misma inyección extraordinaria europea que entonces explicaba por qué "Asuntos Económicos" pesaba tanto dentro del PRTR es, ahora, la que explica por qué esta misma división cae de forma tan visible en la serie agregada según ese impulso extraordinario se agota.

Defensa: la división de menor peso relativo, y el debate de mayor actualidad política de todo el proyecto

[Seguro] "Defensa" es, con diferencia, la división de menor peso absoluto de toda la tabla —apenas 1,0 a 1,2 puntos de PIB entre 2021 y 2024—, pero es también la que ha protagonizado, en el tiempo transcurrido desde que se elaboraron estas cifras, el debate fiscal y geopolítico más intenso y más actual de cuantos hemos recorrido en toda esta serie de tres temas. [Seguro] España destinaba, en 2014, apenas un 0,92% de su PIB a defensa; ese porcentaje subió al 1% en 2020, y a 1,28% en 2024 —la última cifra que la tabla original recogía—. [Seguro] En 2025, según las propias estimaciones de la OTAN, España alcanzó por primera vez en su historia el umbral del 2% del PIB, acordado inicialmente en la Cumbre de Gales de 2014 y cuyo cumplimiento se había ido postergando en sucesivas cumbres posteriores —un salto de más del 43% en un solo año, de 22.693 a 33.123 millones de euros, financiado mediante una inyección adicional de 10.471 millones aprobada expresamente por el Gobierno—. [Seguro] Pero este mismo umbral del 2%, alcanzado con tanto esfuerzo, quedó ya superado por un nuevo objetivo aprobado en la Cumbre de la OTAN de La Haya, celebrada el 24 y 25 de junio de 2025: los 32 Estados miembros de la Alianza se comprometieron, en su Declaración Final, a invertir el 5% del PIB anualmente en necesidades de defensa y seguridad de aquí a 2035 —desglosado en un 3,5% de inversión militar pura y un 1,5% adicional de gasto en seguridad vinculado a la defensa—. [Seguro] España, sin embargo, mantuvo una posición diferenciada dentro de la propia Alianza: el Gobierno defendió, tanto antes como después de la cumbre, que España podía alcanzar las mismas capacidades militares comprometidas con sus aliados destinando únicamente el 2,1% de su PIB, sin necesidad de asumir el objetivo general del 5%. [Probable] Esta posición generó una fricción explícita con la propia Secretaría General de la OTAN —que llegó a declarar, tras la cumbre, "hay un acuerdo sobre no estar de acuerdo" respecto al caso español— y con el Gobierno de Estados Unidos, cuyo Secretario de Estado calificó la posición española de "gran problema" para la Alianza. [Probable] En el plano interno español, la posición del Gobierno generó, a su vez, reacciones divididas: mientras la oposición de centro-derecha criticaba la gestión de la negociación con la OTAN, la propia coalición de Gobierno mostró fisuras visibles —el socio minoritario Sumar expresó resistencias internas al incremento del gasto militar, y formaciones más a la izquierda como Podemos llegaron a acusar al Ejecutivo de mantener una posición insostenible frente al propio texto firmado por España en La Haya, que sí recoge el compromiso general del 5%—. [Suposición] Es, sin necesidad de que este tema tome partido en un debate que sigue sin resolverse en el momento de esta redacción, el ejemplo más vivo y más actual posible de todo lo que hemos construido a lo largo de estos tres temas: una división de la COFOG que apenas pesaba un punto del PIB hace pocos años, convertida hoy en el epicentro de una negociación internacional, de una disputa dentro de la propia coalición de Gobierno, y de una trayectoria de crecimiento presupuestario —del 0,92% de 2014 a un horizonte que varios aliados sitúan, para España, muy por encima del 2,1% que el Gobierno ha defendido— que ningún lector de las cifras estáticas de la tabla, sin este contexto, podría anticipar por sí solo.

Verificación final del Bloque 3 y de todo el Tema 23

Este último epígrafe ha cubierto: la tabla completa de las diez divisiones COFOG con su evolución 2021-2024 y el peso agregado del gasto social; la distribución razonada por subsectores, aplicando con coherencia todo lo ya fundamentado en el Bloque 1 sobre la naturaleza territorial o funcional de cada uno; el declive de Asuntos Económicos explicado por el agotamiento progresivo del PRTR, cerrando el círculo abierto en el Tema 22; y el caso de Defensa, desarrollado con la máxima actualidad disponible y con la neutralidad política que un debate de esta naturaleza exige. No identifico ninguna vertiente pendiente dentro del alcance de este bloque ni de este tema en su conjunto.































Reglamento (UE) nº 549/2013 del Parlamento Europeo y del Consejo, de 21 de mayo de 2013, relativo al Sistema Europeo de Cuentas Nacionales y Regionales (SEC-2010), sectorización institucional.
STONE, R. (1953 y 1968): A System of National Accounts (ya citado en el Tema 22).
ATKINSON, A. B. (2005): The Atkinson Review: Final Report — Measurement of Government Output and Productivity for the National Accounts, Palgrave Macmillan.
ATKINSON, A. B. (2005): "Measurement of UK Government Output and Productivity for the National Accounts", Journal of the Statistical and Social Inquiry Society of Ireland, vol. XXXIV.
Office for National Statistics (Reino Unido): Bean Review of Economic Statistics, 2016, y National Statistician's Independent Review of the Measurement of Public Services Productivity, 2023-2025 (continuidad de los principios del Informe Atkinson).
Naciones Unidas, Comisión Europea, FMI, OCDE, Banco Mundial (2008): Sistema de Cuentas Nacionales 2008 (SCN 2008), revisión de la medición de la producción no de mercado.
- Reglamento (UE) nº 549/2013 (SEC-2010), secuencia de cuentas: producción, distribución y utilización de la renta, capital, financiera, y balance.
- Eurostat: Manual on Government Deficit and Debt — Implementation of ESA 2010, ediciones sucesivas (tratamiento estadístico de titulizaciones, operaciones de venta y arrendamiento posterior, y colaboraciones público-privadas).
- BALZLI, B. (2010): "Greek Debt Crisis: How Goldman Sachs Helped Greece to Mask its True Debt", Der Spiegel, 8 de febrero de 2010.
- DUNBAR, N. (2003): revelación original de la operación de permuta de divisas entre Grecia y Goldman Sachs.
- Eurostat: revisión de las cifras de déficit público griego de 2002, septiembre de 2004 (corrección del 1,2% al 3,7%, posteriormente al 5,2%).
- Tratado de Maastricht (Tratado de la Unión Europea), 1992, criterios de convergencia de déficit y deuda, ya citado en el Tema 21.
- Reglamento (UE) nº 549/2013 (SEC-2010), codificación de subsectores S.1311 a S.1314.
- OCDE: Fiscal Decentralisation Database y Revenue Statistics, ediciones recientes (clasificación de España como "país regional", junto a Colombia, y evolución de la recaudación tributaria autonómica 1995-2021).
- LAGO PEÑAS, S.: "La descentralización tributaria en España: avances significativos, retos pendientes", Cuadernos de Información Económica, FUNCAS.
- Fundación Manuel Giménez Abad: Comparación internacional de los modelos de descentralización fiscal (Alfonso Utrilla de la Hoz), datos de gasto regional en proporción del PIB comparado con Bélgica, Alemania y Austria.
- MORENO, L.: literatura sobre la "federalización sin federalismo" del Estado autonómico español.
- Constitución Española de 1978, art. 2 (unidad de la Nación española y derecho a la autonomía de nacionalidades y regiones).
- IGAE: *Informe trimestral de las Administraciones Públicas*, 4T 2022 (ya citado).
- Gobierno de España: *Actualización del Programa de Estabilidad 2023-2026* (ya citado).
- Comisión Europea (2024): *2024 Ageing Report: Economic and Budgetary Projections for the EU Member States (2022-2070)*, y *Country Fiche for Spain*.
- DE LA FUENTE, Á. (2024): "Notas sobre las proyecciones para España del Ageing Report de 2024", FEDEA, Apuntes 2024/15.
- DOMÉNECH, R. (2024): "El gasto creciente en pensiones públicas", BBVA Research / *Actualidad Económica*, *El Mundo*.
- Real Decreto-ley 2/2023 (ya citado en el Tema 22), mecanismo de activación correctora del Fondo de Reserva vinculado al Informe de Envejecimiento.
- DE LA FUENTE, Á. (2023a): "Observatorio de las CCAA: las cuentas autonómicas en 2022 y entre 2003 y 2022", FEDEA (ya citado).
- IGAE (2023b): *Contabilidad nacional. Serie anual. Operaciones no financieras. Total sector AAPP y sus subsectores* (ya citado).
- ESPING-ANDERSEN, G. (1990): *The Three Worlds of Welfare Capitalism*, Polity Press, Cambridge.
- BROCKHOFF, S., ROSSIGNOL, S. y TAUGOURDEAU, E. (2012): "The Three Worlds of Welfare Capitalism Revisited", sobre la financiación diferenciada de los modelos bismarckiano, beveridgiano y liberal.
- BISMARCK, O. von: legislación social alemana de 1883-1889 (seguro de enfermedad, accidentes e invalidez/vejez), origen del modelo de protección social vinculado a cotizaciones sobre el empleo.
- BEVERIDGE, W.: *Social Insurance and Allied Services* (Informe Beveridge), 1942, origen del modelo de protección social británico de cobertura universal financiada por impuestos generales.
- Reglamento Delegado (UE) nº 877/2013 de la Comisión, de 27 de junio de 2013 (metodología de cálculo de la presión fiscal europea).
- OCDE: *Revenue Statistics*, serie publicada de forma ininterrumpida desde 1965.
- OCDE: *Taxing Wages*, serie anual comparada sobre cuña fiscal del trabajo.
- Comisión Europea (2023a): *Data on Taxation Trends* (ya citado).
- FERNÁNDEZ, R., PONZ, J. y TAGUAS, D. (1995): "Algunas reflexiones sobre la fiscalidad del trabajo..." (ya citado en el Tema 23 original).
- BOSCÁ, J. E., DOMÉNECH, R., FERRI, J. y MÉNDEZ, J. R. (2020): "Financial and fiscal shocks in the great recession and recovery of the Spanish economy", *European Economic Review* (ya citado).
- BOSCÁ, J. E., DOMÉNECH, R. y FERRI, J. (2022): "Cotizaciones sociales, salarios públicos y pacto de rentas", BBVA Research (ya citado).
- FERNÁNDEZ AGUILAR, C. (2022): "¿Ajustar o no ajustar los tramos de IRPF con la inflación?", *A vueltas con la economía*, Universidad Internacional Isabel I.
- Real Decreto-ley de deflactación de tramos del IRPF de 2008 y reforma fiscal de 2015 (Ley 26/2014), como los dos únicos ajustes de la tarifa estatal en dos décadas.
Naciones Unidas, Departamento de Asuntos Económicos y Sociales, División de Estadística (2000): Classifications of Expenditure According to Purpose, Statistical Papers Series M, núm. 84, Nueva York (publicación fundacional conjunta de COFOG, COICOP, COPNI y COPP).
Comisión de Estadística de las Naciones Unidas: aprobación de COFOG-1999, trigésima sesión, marzo de 1999.
Reglamento (CE) nº 113/2002 de la Comisión, de 23 de enero de 2002 (base legal de la aplicación de COFOG en el ámbito de la Unión Europea, dentro del SEC).
Comisión de Estadística de las Naciones Unidas: Decisión 53/125(i), quincuagésimo tercera sesión, 2022 (inicio del proceso de revisión de la COFOG, constitución del equipo de trabajo TT-COFOG bajo el Comité de Expertos de las Naciones Unidas en Clasificaciones Estadísticas Internacionales).
Naciones Unidas, Departamento de Asuntos Económicos, Manual para la clasificación económica y funcional de los gastos del Estado, 1958 (ya citado en el Tema 22, como antecedente genealógico directo de la COFOG).

BAUMOL, W. J. y BOWEN, W. G. (1966): Performing Arts, The Economic Dilemma, Twentieth Century Fund, Nueva York (formulación original de la enfermedad de los costes).
BAUMOL, W. J. (2012): The Cost Disease: Why Computers Get Cheaper and Health Care Doesn't, Yale University Press.
FERRIS, J. S. y WEST, E. G. (1996): "The Cost Disease and Government Growth: Qualifications to Baumol", Public Choice, vol. 89, núm. 1-2.
Comunicación de la Comisión Europea (2022): Marco Temporal de Crisis relativo a las medidas de ayuda estatal destinadas a respaldar la economía a raíz de la agresión de Rusia contra Ucrania, marzo de 2022.
Inversión pública en los procesos de consolidación fiscal: FMI, Fiscal Monitor
Reglamento (UE) 2021/241, de 12 de febrero de 2021 (Mecanismo de Recuperación y Resiliencia, ya citado en el Tema 22).
OTAN: Declaración de la Cumbre de Gales, septiembre de 2014 (objetivo del 2% del PIB en defensa).
OTAN: Declaración Final de la Cumbre de La Haya, 24-25 de junio de 2025 (objetivo del 5% del PIB para 2035).
Ministerio de Defensa de España: cifras de gasto en defensa 2014-2025 según metodología OTAN.
Declaraciones públicas del Gobierno de España, de la Secretaría General de la OTAN y del Departamento de Estado de los Estados Unidos sobre la posición española respecto al objetivo de gasto en defensa, junio-agosto de 2025.










































\clearpage
\section{Conclusión} 


\subsection{Recapitulación}


\subsection{Extensiones y relación con otras partes del Temario}



\subsection{Opinión}

\subsubsection{Idea final (Salida o cierra)}

\clearpage
\pagenumbering{gobble}
\MyBibliography



\MyBibItem{DAWID y GATTI}{Chapter 2 - Agent-Based Macroeconomics}{2018}{https://doi.org/10.1016/bs.hescom.2018.02.006}


% ANEXOS
\clearpage










\end{document}